% Latin 11195 — lecture modernisée du volume entier.
%
% Pass: Opus 5 (claude-opus-5), 2026-09-19 — first pass, unchecked against the
% views by a human, and an interpretation rather than a record. Where this file
% and the transcriptions differ in substance, the transcriptions
% (batch-01.fr.tex, batch-02.fr.tex) are the record.
%
% Sources read: the two batch transcriptions of this volume, in full, twice —
% once for the mathematics, once for the apparatus — and nothing else. No image
% was consulted. No printed edition was consulted, and that matters more here
% than usual, because three printed books stand behind these leaves and all
% three were deliberately left shut: Viète's « De recognitione et emendatione
% æquationum », whose school's notation the Latin treatise writes; Galileo's
% « Discorso intorno alle cose che stanno in su l'acqua », whose estimate of the
% weight of air the Italian treatise quotes at second hand; and Flayder's « De
% arte volandi » (Tübingen 1627), whose title page one leaf of the volume copies
% out. Nothing was taken from any of them, nor from any modern edition or study,
% to fill a gap or settle a word. Where the transcription has a gap, this
% reading says so. Both batches were transcribed on Opus 5; together they cover
% views 1–38, the whole volume.
%
% What the volume turned out to be. NOT one text, and not what the catalogue's
% title says. Two unrelated dossiers bound together, in six hands:
%
%   v. 2 R – 23 R   « De Recognitione Æquationum », Latin, one regular
%                   seventeenth-century cursive, in Viète-school species
%                   notation: the doctrine of how an equation is composed out of
%                   its roots and how it is read back. General doctrine and the
%                   three kinds of side (3 R); quadratics (4 R–5); cubics
%                   (6–11), with « De Determinationibus Æquationum Cubicarum »
%                   (8 R–11); quarto-quadratics (12–20), Caput 9um (21–22),
%                   Caput 10mus Et ultimum (23), stopping in mid-page on « de
%                   iis plura non dicemus ».
%   v. 26 R         « Ars Volandi », a title leaf, a second hand.
%   v. 27 R         Flayder's printed title page copied out by hand, a third
%                   hand. The volume's only date: 1627.
%   v. 28 – 29 R    Latin notes and extracts on human flight, a fourth hand.
%   v. 32 L         A French legend of the letters of the facing figure, a fifth
%                   hand.
%   v. 32 R, 37 R   Two pen drawings of the machine as a winged dragon.
%   v. 33 R – 37 L  « Il volare non è impossibile come fin hora vniuersalmente è
%                   stato creduto », a whole treatise in Italian, a sixth hand.
%
% The volume is therefore UNSORTED, and this reading says so instead of
% inventing a thread between the algebra and the flight dossier. Each dossier
% gets its own run of sections.
%
% NO LETTERS. The correspondence of Roberval, Huygens, Torricelli and Fermat
% that the catalogue's title promises is in neither batch. The résumé says this
% plainly, and so does the closing section, because a reader arriving from the
% catalogue has to be told what is not here.
%
% Notation translated, each with a footnote at its first use. « ∥ », the double
% vertical stroke the transcriptions set \parallel, is the period's sign of
% equality and is modernised to « = » throughout. The homogeneous constants
% named by their dimension — z pl., T Sol., S pl.pl., P Sol., R sol., G pl.,
% B pl., D pl., H pl. — are explained once by Viète's law of homogeneity and
% then dropped, the equations being written monic in the modern way. « Latus »
% is a root; positiva supra a positive root, positiva infra a negative one,
% fictitia a complex one. « Potestas » is the highest power, « affectio » a
% lower-degree term, « antithesis » transposition, « applicare » to divide,
% « ducere » to multiply, « constitutio » the way an equation is put together
% from its roots, « interpretatio » its canonical form, « determinatio » the
% case of repeated roots, « limitatio » the inequality on the coefficients
% outside which the roots cease to be real.
%
% Modern names attached, each where the statement genuinely coincides and each
% footnoted as later than the leaf: the relations between coefficients and roots
% and the elementary symmetric functions; the factor theorem; the discriminant
% of the quadratic and of the cubic; the repeated root as the root of
% gcd(P, P') and the Euclidean algorithm as what the treatise's « successive
% comparisons » perform; for the flight dossier, Archimedes' principle,
% neutral buoyancy and the equilibrium altitude, impulse, the invariance of the
% relative velocity, the principle of virtual work, wing loading and the
% square–cube law.
%
% Every computation the leaves offer was redone, and the recomputations are this
% reading's own: the twelve products of Caput 9; the numerical example of Caput
% 10 and its check on the root 5; the elimination that yields the double root;
% the three maximum lemmas; the two geometrical cases of view 11; the bound the
% treatise claims for the ratio of two solids on view 16. Where the leaf is
% loose, elliptical or wrong, this reading states what is true and footnotes
% what the leaf has — at the T-bound of the cubic, at the sufficiency of the two
% limitations, at the « exceptio regulæ generalis » of Caput 5, at the
% self-subsisting quadratic, at the sign of one product of Caput 9, at the
% missing formula of its seventh proposition, at the sevenfold ratio of view 16,
% at the annotator's « minorem », and, in the Italian treatise, at the weight of
% air, the sphere of fire, the order of falling metals, the whirled cane, and
% the lever's « time ».
%
% Folios. Neither transcription wrote a single \folio{}, deliberately: two ink
% series run on these leaves — a pagination 1–37 in the treatise's own ink, then
% a leaf series 38–62 numbered on rectos only — and they meet at 37/38 with no
% break in the count, while neither is the Library's pale pencil foliation.
% Nothing here is cited by folio. Views only, by \pagerange{}{} at the head of
% each section, grounded in the transcriptions' own \page{} marks.
%
% No date is inferred that the leaves do not give. \dating{} is empty because
% the catalogue's field is empty; no leaf of the volume carries a date; the only
% date anywhere in it is the 1627 inside the hand-copied Flayder title page. The
% two datings this reading does venture — the flight dossier is later than 1627,
% the Italian treatise later than the estimate of the weight of air it quotes —
% are marked in their footnotes as this reading's inferences and nothing more.
% The treatise on equations is anonymous on the leaf, and so is the Italian one;
% « D. Rob. » names Roberval twice, as the author of lemmas the treatise
% BORROWS. This reading attributes neither text to anybody.
\documentclass[11pt,a4paper]{article}
% The preamble shared by every transcription of the mathematicians' manuscripts in Gallica.
%
% It rests on one idea: **uncertainty compounds, it does not resolve.** A
% transcription that smooths over a doubtful word has destroyed the only thing
% it was made for — a reader must be able to tell, at every point, what was on
% the page and what was supplied. The apparatus macros below are therefore the
% critical apparatus, not decoration.
%
% The same commands are recognised by `scripts/render.mjs`, which produces the
% site's reading view, and by `scripts/tei.mjs`. A macro added here must be
% added there too, or rendering fails loudly — which is the wanted behaviour.
%
% Comments here are English, like the rest of the repository. The documents
% this preamble sets are French, and that is deliberate: see the skills.

\ProvidesPackage{germain}[2026/09/15 Transcriptions of mathematicians' manuscripts in Gallica]

\usepackage{amsmath,amssymb,amsthm}
\usepackage{mathrsfs}
% Kept from the parent preamble so the renderer's diagram support stays
% available; these manuscripts draw no commutative diagrams, and nothing here uses it.
\usepackage{tikz-cd}
\usepackage[english,french]{babel}
\usepackage{fontspec}

% --- Greek ------------------------------------------------------------------
% The text face stays fontspec's default, Latin Modern, which has no Greek:
% XeTeX drops every glyph it lacks with a line in the log and nothing on the
% page, and the Greek words the manuscripts quote vanished from the PDFs. So
% the Greek and Greek Extended blocks get a XeTeX character class of their own,
% and entering or leaving a run of it switches the family to CMU Serif — the
% same Computer Modern design, with polytonic Greek — and back. Only the
% family changes: a Greek word in \textit or \textbf keeps its shape and series.
% The transitions to and from every other class carry the tokens that class
% has towards an ordinary letter, so French spacing before « ; » or inside
% « » still holds after a Greek word. No grouping, so no brace or \endgroup
% can come out unbalanced; maths is left alone, since XeTeX inserts nothing
% there. Kept identical in `exercices.sty`.
\makeatletter
\newfontfamily\germainGreekfont{cmun}[
  Extension=.otf, NFSSFamily=germaingreek,
  UprightFont=*rm, ItalicFont=*ti, SlantedFont=*sl,
  BoldFont=*bx, BoldItalicFont=*bi, BoldSlantedFont=*bl]
\newXeTeXintercharclass\germain@greek
\def\germain@greekrange#1#2{%
  \@tempcnta=#1\relax
  \loop
    \XeTeXcharclass\@tempcnta=\germain@greek
    \ifnum\@tempcnta<#2\advance\@tempcnta\@ne\repeat}
\germain@greekrange{"0370}{"03FF}
\germain@greekrange{"1F00}{"1FFF}
\def\germain@greekprev{\rmdefault}
\def\germain@greekin{%
  \edef\germain@greekprev{\f@family}\fontfamily{germaingreek}\selectfont}
\def\germain@greekout{\fontfamily{\germain@greekprev}\selectfont}
\def\germain@greekjoin{%
  \XeTeXinterchartokenstate=\@ne
  \@tempcnta=\z@
  \loop
    \ifnum\@tempcnta=\germain@greek\else
      \XeTeXinterchartoks\germain@greek\@tempcnta=\expandafter{%
        \expandafter\germain@greekout\the\XeTeXinterchartoks\z@\@tempcnta}%
      \XeTeXinterchartoks\@tempcnta\germain@greek=\expandafter{%
        \the\XeTeXinterchartoks\@tempcnta\z@\germain@greekin}%
    \fi
    \ifnum\@tempcnta<4095\advance\@tempcnta\@ne\repeat}
% babel-french sets its own classes, and saves and restores the interchar
% state, each time a language is selected: join again after it does.
\addto\extrasfrench{\germain@greekjoin}
\addto\extrasenglish{\germain@greekjoin}
\AtBeginDocument{\germain@greekjoin}
\makeatother

\usepackage{geometry}
\geometry{a4paper,margin=25mm,marginparwidth=28mm}
\usepackage{marginnote}
\usepackage{xcolor}
\usepackage[normalem]{ulem}
\setlength{\emergencystretch}{3em}
\usepackage{hyperref}
\hypersetup{colorlinks=true,linkcolor=black,urlcolor=black,pdfborder={0 0 0}}

% --- Batch metadata ---------------------------------------------------------
% Taken as they stand by the rendering script, which shows them at the head of
% the reading view. They fix what the file claims to cover. The macro names are
% the parent project's, so that the scripts stay shared; \folder here means the
% volume's slug — `fr-9115` — and \pages counts Gallica views.

\newcommand{\folder}[1]{\gdef\grothFolder{#1}}
\newcommand{\batch}[1]{\gdef\grothBatch{#1}}
\newcommand{\pages}[2]{\gdef\grothFirst{#1}\gdef\grothLast{#2}}
\newcommand{\foldertitle}[1]{\gdef\grothFoldertitle{#1}}

% The holder's own shelfmark, printed as they write it: « Français 9115 ».
\newcommand{\shelfmark}[1]{\gdef\grothShelfmark{#1}}

% The dating, copied verbatim from the holder's catalogue — never inferred
% here. « XIXe siècle » is the BnF's; a pass that dates a leaf from its content
% says so in a \note{}, as its own inference.
\newcommand{\dating}[1]{\gdef\grothDating{#1}}

% The status notice, printed in the title block and repeated as a diagonal
% watermark on every page of the PDF. The manuscripts are in the public
% domain; what this line declares is not a right but a quality — an
% unverified first pass by a machine. The skills write the line; a file
% without it gets no watermark, so leaving it out is visible in the output.
\newcommand{\watermark}[1]{\gdef\grothWatermark{#1}}

\gdef\grothFolder{?}\gdef\grothBatch{}\gdef\grothFirst{?}\gdef\grothLast{?}%
\gdef\grothFoldertitle{}\gdef\grothShelfmark{}\gdef\grothDating{}\gdef\grothWatermark{}

% --- The résumé -------------------------------------------------------------
\newenvironment{resume}
  {\small\setlength{\parskip}{0.45em}}
  {\par\medskip\hrule\medskip}

% --- Keywords ---------------------------------------------------------------
% In English, closing the modernised reading's résumé: the modern vocabulary
% under which to search for what the leaves construct. The manifest extracts
% them to tag the volume — this line is the single source of the tags.
\newcommand{\keywords}[1]{\par\smallskip\noindent
  {\footnotesize\textsc{Keywords} --- \textit{#1}}\par}

% --- The critical apparatus -------------------------------------------------

% The Gallica view number, in the margin. This is the anchor that lets a
% disagreement over a formula be settled by looking at *one* image rather than
% a volume of 750. The site uses it to turn the facsimile as the transcription
% is scrolled. A view is one scanned image: usually one side of a leaf.
\newcommand{\page}[1]{%
  \marginnote{\footnotesize\color{gray!70}\textsf{v.\,#1}}%
  \label{page:#1}\ignorespaces}

% The BnF's pencilled foliation, where the leaf carries it and a pass can read
% it: « 198r ». This is what the literature cites — « ff. 198r–208v » — and
% the one line that turns a citation into a view. Recorded, never inferred:
% a folio a pass cannot read is not written.
\newcommand{\folio}[1]{%
  \marginnote{\footnotesize\color{gray!50}\textsf{f.\,#1}}\ignorespaces}

% The modernised reading's coarser counterpart to \page{}: which views a
% section is drawn from.
\newcommand{\pagerange}[2]{%
  \marginnote{\footnotesize\color{gray!70}\textsf{v.\,#1--#2}}%
  \ignorespaces}

% Illegible: never guessed.
\newcommand{\ill}{\textcolor{red!60!black}{[\,\ldots\,]}}

% A doubtful reading: the word is given, and flagged as uncertain.
\newcommand{\uncertain}[1]{\uline{#1}}

% An editorial addition — a missing letter, an implied word.
\newcommand{\add}[1]{\textcolor{blue!50!black}{[#1]}}

% Struck out by the writer. What was crossed out is part of the document.
\newcommand{\struck}[1]{\sout{#1}}

% The transcriber's note, on the state of the page or the mathematics.
\newcommand{\note}[1]{\par\smallskip{\small\color{gray!60!black}\textit{[#1]}}\par\smallskip}

% A marginal note of hers, distinct from the body of the text.
\newcommand{\marginal}[1]{\par\smallskip\noindent\hspace{1em}%
  {\small\color{gray!60!black}#1}\par\smallskip}

% --- Title ------------------------------------------------------------------
% The title block is French, like the documents themselves.

\AddToHook{shipout/background}{%
  \ifx\grothWatermark\empty\else
    \begin{tikzpicture}[remember picture, overlay]
      \node[rotate=52, scale=3.4, align=center, text=gray!18,
            font=\sffamily\bfseries]
        at (current page.center) {\grothWatermark};
    \end{tikzpicture}%
  \fi}

\AtBeginDocument{%
  \begin{center}
    {\large\bfseries Manuscrits de mathématiciens (Gallica) ---
      \ifx\grothShelfmark\empty\grothFolder\else\grothShelfmark\fi}\\[2pt]
    {\itshape\grothFoldertitle}\\[4pt]
    {\small vues \grothFirst--\grothLast
      \ifx\grothBatch\empty\else\ (lot \grothBatch)\fi}\\[2pt]
    \ifx\grothDating\empty\else
      {\small\color{gray!60!black}Datation du catalogue : \grothDating}\\[2pt]
    \fi
    \ifx\grothWatermark\empty\else
      {\small\bfseries\color{red!45!gray}\grothWatermark}\\[2pt]
    \fi
    {\footnotesize\color{gray!60!black}Transcription établie d'après les images IIIF de Gallica.
     Source gallica.bnf.fr / Bibliothèque nationale de France.\\
     Les lectures incertaines sont soulignées, les illisibles notées [\,\ldots\,],
     les ajouts entre crochets ; \textsf{v.} numérote les vues de Gallica,
     \textsf{f.} la foliotation au crayon de la BnF.}
  \end{center}
  \vspace{1em}}

\endinput


\folder{latin-11195}
\shelfmark{Latin 11195}
\pages{1}{38}
\dating{}
\watermark{Lecture automatique\\interprétation non vérifiée}
\foldertitle{Mélanges de physique et de mathématiques, comprenant des opuscules et des lettres de P. de Roberval, Huggens, Torricelli, Firmat et Fr. Herman Flayder. Latin 11195 — lecture modernisée du volume entier}

\begin{document}

\section*{Résumé}

\begin{resume}
Ce volume n'est pas un livre : c'est une reliure. Deux dossiers qui n'ont rien à
voir l'un avec l'autre y ont été cousus ensemble, et six mains au moins s'y
partagent trente-huit vues. Le premier dossier est un traité d'algèbre en latin,
anonyme, sur la manière dont une équation est fabriquée à partir de ses racines
et dont on peut, en la regardant, remonter à celles-ci. Le second est un dossier
sur le vol humain : un feuillet de titre, la page de titre d'un livre imprimé
recopiée à la main, des notes latines rassemblant tout ce que leur auteur avait
lu sur des hommes qui avaient volé ou cru voler, deux dessins à la plume d'une
machine en forme de dragon, la légende française des lettres de l'un d'eux, et
un traité entier en italien qui soutient que voler n'est pas impossible et
décrit la machine que son auteur dit avoir inventée. Entre les deux dossiers,
aucun lien : ni la main, ni l'encre, ni le sujet, ni la moindre phrase de
renvoi. Cette lecture ne leur en invente pas.

Il faut dire tout de suite ce que le volume ne contient pas, parce que le titre
du catalogue induit en erreur. Il annonce des opuscules et des lettres de
Roberval, de Huygens, de Torricelli et de Fermat : il n'y a ici aucune lettre,
d'aucun d'eux. Roberval n'y paraît que sous la forme « D. Rob. », deux fois,
comme l'auteur de lemmes que le traité d'algèbre emprunte, avec une note
marginale en français qui renvoie le lecteur à « son traité de mechanique des
plans inclinez ». C'est la preuve que le traité cite Roberval, non qu'il soit de
lui. Aucun feuillet ne porte de date ; la seule date du volume, 1627, est à
l'intérieur de la page de titre recopiée.

Le traité d'algèbre pose un problème simple à énoncer et difficile à tenir. On
sait fabriquer une équation : on prend des racines, on écrit pour chacune un
facteur, on multiplie, et l'on obtient une équation dont les coefficients sont
des expressions de ces racines. Peut-on faire le chemin inverse — lire, dans les
coefficients d'une équation donnée, combien de racines elle a, de quel signe, et
comment les atteindre ? Le traité appelle cela la \emph{reconnaissance} de
l'équation, et il le dit dans sa première phrase : « on y cherche les moyens par
lesquels les diverses constitutions des équations se reconnaissent » ; car « une
équation contient en elle toutes les racines dont elle a été composée, et c'est
donc de toutes celles-là qu'elle est explicable ». L'idée qui porte tout
l'ouvrage est un inventaire : on énumère toutes les manières d'assembler une
équation d'un degré donné à partir de racines de signes donnés, on développe
chaque produit, on nomme les coefficients qui en sortent, et l'on obtient un
catalogue de formes canoniques dans lequel il ne reste plus, devant une équation
inconnue, qu'à retrouver la sienne. Quatre-vingt-sept propositions plus tard, le
catalogue est complet pour les degrés deux, trois et quatre, et le traité
s'arrête au milieu d'une page sur « nous n'en dirons pas plus ».

Ce que ce catalogue a trouvé, et que l'on peut aujourd'hui nommer, ce sont les
relations entre les coefficients d'une équation et les fonctions symétriques
élémentaires de ses racines : le traité les énonce en toutes lettres et pour
tous les degrés. Ce qui l'empêche de les tenir, c'est qu'il ne compte pas les
racines complexes comme des racines. Il les appelle \emph{fictitia}, il les
reconnaît fort bien — il sait dire à quelle condition une quadratique en a, et
il construit exprès une quarto-quadratique qui en a deux — mais il leur refuse
le statut de racine, et sa belle règle générale se trouve alors avoir des
exceptions. Cette lecture montre, sur les exemples du feuillet lui-même, que ces
exceptions disparaissent toutes dès qu'on compte les côtés fictifs comme des
racines. Le reste du traité est un travail de première qualité sur les
inégalités : les « limitations » qu'il établit, degré par degré, en résolvant à
chaque fois un problème de maximum, sont exactement les conditions que nous
écrivons aujourd'hui sur le discriminant. Et un béquet épinglé sur un feuillet
blanc pose, en trois lignes, la seule question à laquelle la doctrine ne sait
pas répondre — comment distinguer deux formes qui ont les mêmes signes ? — dont
la réponse moderne est précisément le signe du discriminant.

Le traité italien, lui, veut prouver que voler est possible et le prouver en
mécanicien. Son auteur raconte pourquoi : « depuis mes premières années j'ai
appliqué mon esprit à cette invention, trouvant étrange que l'esprit humain ne
fût pas capable d'atteindre par son savoir aux choses qui ont été accordées par
nature à d'autres animaux inférieurs à nous » ; il a essayé, tout fut inutile,
puis « voici dix ans écoulés » qu'il a trouvé, « à partir de principes assez
communs », ce qu'il va exposer. Son raisonnement va du poids de l'air à la
flottaison, de la flottaison à la percussion, et de la percussion au battement
des ailes : rien de plus léger que l'air n'étant à notre portée, il faut se
soutenir en frappant l'air, comme les oiseaux, et il divise la percussion en
trois espèces selon que le corps frappé est immobile, vient à la rencontre du
coup, ou le fuit. Puis, sur un levier, il démontre qu'on ne peut gagner à la
fois la force et le temps, et il en tire, honnêtement, la limite de sa propre
machine : elle ne pourra pas porter plus de deux hommes. Beaucoup de ce qu'il
écrit est juste — la proposition d'Archimède qu'il cite, les pesanteurs
spécifiques de cinq métaux à cinq pour cent près, la loi du levier, le rôle de
la vitesse relative dans un choc, le repliement des ailes à la montée ; et la
machine ne pouvait pas voler, pour une raison que ses propres nombres
permettent d'établir et qu'il n'a pas vue, qui est la puissance.

Où cela mène. Du côté de l'algèbre : à la théorie des équations, aux relations
entre coefficients et racines, aux fonctions symétriques élémentaires, au
discriminant, au plus grand commun diviseur de $P$ et $P'$ pour les racines
multiples, et à l'histoire longue de l'admission des nombres négatifs puis
complexes comme racines de plein droit. Du côté du vol : à la poussée
d'Archimède et à la flottaison, à l'altitude d'équilibre d'un ballon, à la
notion d'impulsion, au principe des travaux virtuels, à la charge alaire et à la
loi carré-cube, et à l'aérodynamique du vol battu.

\keywords{theory of equations, Viète's species notation, law of homogeneity,
elementary symmetric functions, relations between roots and coefficients,
negative roots, complex roots, discriminant, repeated roots, greatest common
divisor of a polynomial and its derivative, cubic equations, quartic equations,
cubic extremum problems, Roberval, history of flight, ornithopter, buoyancy,
density of air, Archimedes' principle, specific gravity, percussion, impulse,
law of the lever, virtual velocities, wing loading, square-cube law, Flayder,
Galileo, Sagredo, manuscript miscellany}
\end{resume}

\section*{Ce que ce volume contient}

\pagerange{2}{38}
Trente-huit vues, chacune montrant le volume ouvert, que les deux
transcriptions ont coupées au pli et lues comme deux pages. Neuf vues sont
entièrement blanches et ne sont pas transcrites — la vue 1, les vues 24 et 25,
30, 31 et 38, ainsi que les pages de droite des vues 14 et 15 ; le trou dans la
numérotation est le témoignage. Tout le reste se partage en deux dossiers sans
rapport l'un avec l'autre.

\begin{itemize}
  \item \textbf{Vues 2 R à 23 R — un traité d'algèbre en latin}, anonyme, sans
    titre d'auteur, intitulé sur son feuillet de titre « De Recognitione
    Æquationum », d'une seule cursive régulière que les transcriptions
    décrivent comme une écriture du XVII\textsuperscript{e} siècle. Il est
    complet de son début et s'arrête de lui-même au milieu de la vue
    23 R.\footnote{C'est une description de l'écriture faite par la
    transcription, non une datation du texte : aucun feuillet ne porte de date,
    et cette lecture n'en infère aucune. Voir la dernière section.}
  \item \textbf{Vues 26 R à 37 R — un dossier sur l'art de voler}, en cinq
    pièces et quatre écritures au moins : un feuillet de titre portant les deux
    mots « Ars Volandi » (26 R) ; la page de titre d'un livre imprimé, celui de
    Flayder, recopiée à la main et centrée sur le feuillet (27 R) ; des notes et
    extraits latins sur le vol, sans titre, d'une cursive fine et très rapide,
    qui s'interrompent au milieu de la vue 29 R ; la légende française des
    lettres d'une figure (32 L) ; et un traité entier en italien, « Il volare
    non è impossibile come fin hora vniuersalmente è stato creduto » (33 R à
    37 L). Deux dessins à la plume de la même machine occupent les vues 32 R et
    37 R.
\end{itemize}

Les deux dossiers ne se touchent nulle part. Aucune main n'est commune, aucune
phrase de l'un ne renvoie à l'autre, et les sujets n'ont rien à voir. Le second
dossier est du reste séparé du premier par la lacune que signale un feuillet
moderne posé sur la vue 25 : « folios blancs 38 à 49 ». Cette lecture donne donc
à chaque dossier sa propre suite de sections et ne cherche pas de fil entre
eux.\footnote{Sur le compte des mains : les transcriptions distinguent la main
du traité d'algèbre, puis, dans le dossier du vol, celle qui a recopié la page
de titre de Flayder (« une cursive fine et régulière, différente de celle du
traité des équations »), celle des notes latines (« différente de celle du
traité des équations comme de la belle ronde qui a copié la page de titre »),
celle de la légende française (« une troisième main, française, plus ronde, plus
grande ») et celle du traité italien (« une quatrième main »). La main des deux
mots « Ars Volandi » de la vue 26 R est dite seulement différente de celle du
traité ; savoir si elle est ou non celle de la page de titre recopiée n'est pas
tranché par les transcriptions, et ne l'est pas ici. Le compte de six mains
retenu dans cette lecture suppose qu'elle en est une ; cinq est un minimum
certain.}

Pourquoi les feuillets ne sont cités ici que par les vues de Gallica. Deux
séries de chiffres courent sur ces pages, toutes deux à l'encre. La première est
la pagination du traité, de 1 (vue 3 R) à 37 (vue 23 R), un chiffre par page, au
coin extérieur supérieur, de l'encre du texte. La seconde ne numérote que les
pages de droite, un chiffre par feuillet, en chiffres plus grands et plus
ornés : 38 à la vue 24, puis 50 à 62 sur les vues 26 à 38, sans lacune. Les deux
séries se rejoignent à 37 et 38 sans rupture dans le compte, ce qu'aucune des
deux transcriptions n'a su expliquer ; et ni l'une ni l'autre n'est le crayon
pâle d'une foliotation de bibliothèque. Aucune n'a donc écrit de foliotation, et
cette lecture n'en écrit pas : elle ne cite que des vues.\footnote{Les autres
chiffres des feuillets, tous relevés dans les notes des transcriptions, sont : le
« 171 » du feuillet de titre, rogné par le couteau du relieur et d'une autre
main ; les anciennes cotes « Suppl.\textsuperscript{t} lat. 218. A. » sur la
garde et « Suppl.\textsuperscript{t} Lat. n\textsuperscript{o} 218. » à la vue 3 ;
un « 26 » au dos du feuillet de titre ; une série de petits chiffres courants de
1 à 29 dans la marge du traité, d'une seconde plume plus fine, et une seconde
série de la même sorte, « 29\textsuperscript{b} » à 37, dans les marges des vues
21 à 23 ; les numéros de proposition ; et, aux vues 19 et 20, un « 33 » qui est
la pagination d'un feuillet postérieur, vue par avance sur le bord saillant.}

Une seconde plume, plus fine et plus pâle, a travaillé dans les marges du
traité : elle y a recopié la forme canonique de chaque constitution en regard de
celle-ci, numéroté ses propres remarques de 1 à 29, posé la vérification chiffrée
d'un lemme à la vue 9 R, et proposé une correction à la vue 11 R. Ses
interventions sont signalées ici chaque fois qu'elles portent sur la substance ;
l'une d'elles est fausse et on le dira.

\section*{Le traité « De Recognitione Æquationum » : ce qu'il se propose}

\pagerange{2}{3}
Le titre est trois mots au milieu du tiers supérieur d'un feuillet, et le traité
commence à la vue suivante en disant exactement ce qu'il fait :

\begin{quote}
« Ideo is tractatus dicitur de æquationum recognitione quia in ipso
investigantur modi quibus æquationum variæ constitutiones dignoscantur, seu per
quos explicentur. »\footnote{« Ce traité est dit de la reconnaissance des
équations parce qu'on y cherche les moyens par lesquels les diverses
constitutions des équations se distinguent, ou par lesquels elles se
développent. » Les citations latines sont données ici dans l'orthographe et avec
les abréviations développées ; le feuillet écrit « d\textsuperscript{r} » pour
« dicitur », « recognitiōe » pour « recognitione », « modj » pour « modi ». La
transcription conserve toutes ces formes telles quelles.}
\end{quote}

La raison pour laquelle il faut un tel art est donnée aussitôt : les opérations
de l'analyse — multiplications, divisions — cachent les équations « sous des
formules diverses, en sorte que leurs racines, ou côtés, enveloppées dans ces
détours, se livrent difficilement ». Reconnaître une équation, c'est donc
défaire ce qu'une multiplication a fait. Et le principe qui rend l'entreprise
possible est énoncé dans la phrase suivante, qui est la thèse de tout
l'ouvrage :

\begin{quote}
« Æquatio enim in se habet eas omnes radices ex quibus fuit composita, ac
proinde de iisdem omnibus est explicabilis. »\footnote{« Car une équation
contient en elle toutes les racines dont elle a été composée, et c'est donc de
toutes celles-là qu'elle est explicable. » En termes modernes : si
$P(A) = (A - B_1)\cdots(A - B_n)$, alors $B_1, \dots, B_n$ sont racines de $P$,
et ce sont les seules. C'est ce qu'on appelle aujourd'hui le théorème du facteur
et sa réciproque ; le nom est postérieur au feuillet, et la proposition y est
prise comme évidente, non démontrée.}
\end{quote}

De là le traité tire un corollaire dont il fait son plan : plus le degré est
élevé, plus une équation est « impliquée », et de plus de côtés elle est
explicable, parce qu'elle peut naître de plus de façons. Une quadratique ne peut
venir que du produit de deux latérales ou subsister par elle-même ; une cubique
peut venir de trois latérales, ou d'une quadratique et d'une latérale, ou
subsister par elle-même ; et ainsi de suite. Le traité ajoute une phrase sur ses
prédécesseurs qu'il faut citer parce qu'elle dit ce qu'il croit apporter : « les
anciens analystes, enveloppés d'épaisses ténèbres, étaient là-dessus tout à fait
aveugles, car ils pouvaient à peine exhiber un seul côté ». L'élégance de la
méthode, dit-il, tient à ceci qu'elle exhibe \emph{tous} les côtés dont une
équation quelconque a été composée.\footnote{« Qua in re veteres analystæ densis
offusi tenebris plane cæcutiebant, vix enim unicum poterant latus exhibere,
atque inde etiam patet huius methodi elegantia, quæ in eo, ut diximus, est
præcipua ut omnia omnino latera ex quibus quævis æquatio composita fuerit
exhibeat. » Le feuillet ne nomme aucun de ces anciens et cette lecture n'en
nomme pas davantage. L'ambition annoncée — toutes les racines, non la seule
racine positive — est celle-là même que le traité ne tiendra pas jusqu'au bout,
puisqu'il laisse hors du compte les racines complexes.}

Le plan est donc un inventaire. On prend un degré, on énumère les manières
d'assembler une équation de ce degré, on écrit le produit, on nomme ses
coefficients, et l'on obtient une forme canonique. Chaque façon d'assembler est
un \emph{caput}, chaque configuration de grandeurs à l'intérieur d'un caput est
une \emph{propositio}. L'ouvrage compte ainsi, sur les vingt-deux feuillets
écrits, quatre chapitres et cinq propositions pour les quadratiques, six
chapitres et vingt propositions pour les cubiques, dix chapitres et une
soixantaine de propositions pour les quarto-quadratiques, plus une partie
entière sur les « déterminations » des cubiques. Deux traités compagnons sont
cités et ne sont pas dans le volume : un « tractatus de locis » où les
changements de signes seront mieux expliqués, et un « tractatus peculiaris de
generibus » auquel le dernier chapitre renvoie les déterminations des
quarto-quadratiques, ce qui est la raison pour laquelle le traité
s'arrête.\footnote{« quæ facilius percipientur in tractatu de locis » (vue 4 L)
et « quia facile haberi possunt vel ex iis quæ diximus, vel ex tractatu
peculiari de generibus, ideo, de iis plura non dicemus » (vue 23 R, dernière
phrase). Ces deux renvois montrent que le traité faisait partie d'un ensemble
plus large ; rien dans le volume ne dit lequel, et cette lecture n'en conclut
rien.}

\section*{La notation du feuillet, et sa traduction}

\pagerange{3}{5}
Rien dans ce traité ne peut être lu sans traduire d'abord trois choses : le
signe d'égalité, la loi d'homogénéité qui gouverne les noms des constantes, et
le vocabulaire des racines. Ce sont les trois traductions dont dépend tout le
reste de cette lecture.

\subsection*{Le signe d'égalité}

Le feuillet n'écrit jamais « $=$ ». Il écrit deux barres verticales parallèles,
« $\parallel$ », et cette double barre est son signe d'égalité. Toutes les
équations de cette lecture sont écrites avec « $=$ » ; c'est la première et la
plus fréquente des traductions faites ici.\footnote{Les deux transcriptions ont
choisi indépendamment de rendre ce signe par \texttt{\textbackslash parallel}
et de ne pas le moderniser, au motif que le tracé du feuillet est celui de deux
simples hampes verticales et que rien sur la page n'autorise un « $=$ ». Ainsi
la première équation du traité, à la vue 4 L, s'y lit
« $4 - 1\,\mathrm{R} \parallel 0$ » et se lit ici $4 - R = 0$. Le même signe sert
une fois de signe de proportion, dans la vérification chiffrée du lemme de la
vue 9 R.}

Une convention d'écriture va avec : le traité met toujours tous les termes du
même côté et égale le tout à zéro. Il le dit et il dit pourquoi : ainsi « les
affirmés seront égaux aux niés », c'est-à-dire que la somme des termes précédés
de $+$ égale celle des termes précédés de $-$, et l'on a de quoi vérifier une
racine par un simple calcul de deux colonnes. C'est une commodité de calcul, et
le traité s'en sert effectivement à la dernière page.\footnote{« æquatio […] in
hac methodo sic est constituenda, ut quantitas ignota sit ex eadem parte cum
cognita, et utraque pars æquationis sic in unum redacta æquetur nihilo […] Et
tunc affirmata æqualia erunt negatis » (vue 4 L). L'usage qui en est fait est à
la vue 23 R, sur l'équation numérique du dernier chapitre.}

\subsection*{L'homogénéité, et pourquoi les constantes portent un nom de dimension}

Les constantes du traité ne s'appellent pas $a$, $b$, $c$. Elles s'appellent
« $z\,\mathrm{pl.}$ », « $T\,\mathrm{Sol.}$ », « $S\,\mathrm{pl.pl.}$ »,
« $P\,\mathrm{Sol.} $», « $R\,\mathrm{sol.}$ », « $G\,\mathrm{pl.}$ ». Le mot qui
suit la lettre n'abrège pas la lettre : il donne sa \emph{dimension}.
« $z\,\mathrm{pl.}$ » est $z$ \emph{planum}, un plan ; « $T\,\mathrm{Sol.}$ » est
$T$ \emph{solidum}, un solide ; « $S\,\mathrm{pl.pl.}$ » est $S$
\emph{planoplanum}, un plano-plan, c'est-à-dire un produit de quatre longueurs.

La raison en est la loi d'homogénéité de l'école de Viète : on ne peut ajouter
ni égaler que des grandeurs de même dimension. Une longueur est une longueur, le
produit de deux longueurs est un plan, celui de trois un solide, celui de quatre
un plano-plan ; et une somme dont les termes ne sont pas tous de même espèce
n'est pas fausse, elle n'a pas de sens. Il s'ensuit que dans une équation dont
la plus haute puissance est $A^2$, tout terme doit être un plan : le terme
constant n'est donc pas un nombre mais un plan, et il faut l'appeler
$z\,\mathrm{pl.}$ ; le coefficient de $A$, lui, doit être une longueur, et on
l'appelle $x$, sans mot de dimension, parce qu'une longueur n'en demande pas.
Dans une cubique, tout terme doit être un solide : d'où $T\,\mathrm{Sol.}$ seul,
$z\,\mathrm{pl.}\,A$ (plan fois longueur), $xA^2$ (longueur fois plan), $A^3$.
Dans une quarto-quadratique, tout terme est un plano-plan : d'où
$S\,\mathrm{pl.pl.}$, $T\,\mathrm{Sol.}\,A$, $z\,\mathrm{pl.}\,A^2$, $xA^3$,
$A^4$. La règle tient en une ligne : \emph{la dimension du coefficient et le
degré de la puissance qu'il affecte s'ajoutent toujours pour faire le degré de
l'équation.} C'est cela, et rien d'autre, que disent les mots « pl. », « Sol. »
et « pl.pl. ».

L'algèbre moderne a abandonné cette exigence : nous traitons les coefficients
comme des nombres purs, et nous écrivons la quarto-quadratique du premier
chapitre
\[ A^4 - x A^3 + z A^2 - T A + S = 0. \]
C'est ainsi que cette lecture écrira désormais toutes les équations du traité,
les mots de dimension une fois expliqués étant laissés de côté. Le prix de ce
gain est réel et il faut le nommer : chez Viète, chaque équation est un énoncé
sur des figures, et la loi d'homogénéité est la garantie qu'elle en est un. En
la quittant, on perd la raison d'être des noms, et l'on perd du même coup le
contrôle dimensionnel qui, sur ces feuillets, sert plus d'une fois de garde-fou —
c'est lui qui permet de voir tout de suite qu'une phrase comme « un solide dont
la base est $\tfrac23 z\,\mathrm{pl.}$ et la hauteur le tiers du même plan » est
elliptique, une hauteur ne pouvant pas être le tiers d'un plan.\footnote{Cette
phrase est réellement sur le feuillet, à la vue 9 R ; on verra plus bas ce
qu'elle veut dire et pourquoi elle est correcte une fois l'ellipse levée.}

Le reste des lettres suit l'usage de l'école : $A$ est l'inconnue, $B$, $C$,
$D$, $F$ les côtés donnés, $x$ le coefficient du terme immédiatement inférieur
à la plus haute puissance. Le traité ne se sert de la notation cossique des
degrés qu'une seule fois, à la dernière page, quand il passe aux nombres : $a$
pour le côté, $q$ pour le carré, $c$ pour le cube, $qq$ pour le
quarré-quarré.\footnote{« $100 - 105\,a + 7\,q - 3\,c + 1\,qq \parallel 0$ »
(vue 23 R), soit $A^4 - 3A^3 + 7A^2 - 105A + 100 = 0$. C'est la seule équation
entièrement numérique du traité.}

\subsection*{Les côtés : au-dessus, au-dessous, feints}

« Latus », un côté, est ce que nous appelons une racine — le mot est
géométrique, le côté d'un carré ou d'un cube étant ce dont on part pour faire la
puissance. Le traité en distingue trois espèces, et c'est ici que se joue toute
la portée de l'ouvrage.

\begin{itemize}
  \item \textbf{Positiva supra} : les côtés « qui expriment quelque chose outre
    ou au-dessus de rien ». Ce sont nos racines positives. « Si $+A$ est positif
    au-dessus et que sa valeur soit 4, on entendra que la valeur de $A$ est
    quatre outre ou au-dessus de rien. »
  \item \textbf{Positiva infra} : les côtés « conçus comme moindres que rien ».
    Ce sont nos racines négatives. Le traité sait que la notion est heurtante et
    la défend : « bien qu'on les conçoive comme moindres que rien, et que cette
    conception paraisse tout à fait feinte — car la nature ne produit rien
    au-dessous de rien —, ils ne sont pourtant pas feints au plus haut point,
    puisque les signes changés ils deviennent positifs au-dessus ». Et il donne
    la raison de calcul pour laquelle on en a besoin : quand on compare la ligne
    inconnue à une ligne connue, il arrive qu'on n'en prenne qu'une partie, et
    que cette partie soit à retrancher ; on parvient tout de même à la
    connaissance de la ligne, et l'on restitue à la fin ce qu'on avait
    ôté.\footnote{Vue 4 L. Le raisonnement est exactement celui qui justifie,
    aujourd'hui, de travailler sur une droite orientée : la partie
    « au-dessous » est un déplacement de signe contraire. Ce que le feuillet
    n'a pas, c'est le nombre négatif lui-même comme objet ; il a la racine
    négative comme \emph{situation}, décrite par un jeu de signes.}
  \item \textbf{Fictitia} : les côtés qui « ne contiennent rien de vrai en
    eux ». Ce sont nos racines complexes, et le traité dit d'eux la chose
    juste : « si ces feints servent seuls à l'explication d'une équation, ils
    dénoteront l'impossibilité de la question proposée ».\footnote{Vue 3 R. En
    termes modernes : si toutes les racines d'une équation sont complexes non
    réelles, le problème géométrique dont elle est issue n'a pas de solution. La
    proposition est vraie. Le mot « fictitia » est traduit ici par
    « complexe » plutôt que par « imaginaire », parce que le traité ne leur
    donne jamais de forme et ne calcule jamais avec eux ; il ne les connaît que
    par l'échec de la réduction.}
\end{itemize}

Le point décisif est que le traité n'écrit \emph{jamais} une racine négative
comme un nombre négatif. Il écrit le facteur $B + A = 0$, et il dit que l'unique
côté de cette latérale est « $+A$ au-dessous », lequel « par équivalence est le
même que $-A$ au-dessus ». Traduit, cela veut dire $A = -B$, et le facteur est
$A + B$. Ainsi, dans tout ce qui suit :
\begin{itemize}
  \item un facteur écrit $B - A = 0$ apporte une racine positive $A = B$, et
    contribue au produit le facteur $(A - B)$ ;
  \item un facteur écrit $B + A = 0$ apporte une racine négative $A = -B$, et
    contribue le facteur $(A + B)$.
\end{itemize}
Les signes de la forme canonique se déduisent alors immédiatement : le signe du
terme de plus haut degré est $(-1)^k$, où $k$ est le nombre de côtés
« au-dessus », parce que chaque facteur $B - A$ apporte un $-A$ et chaque
facteur $B + A$ un $+A$. C'est ce qui explique pourquoi les chapitres à nombre
pair de racines positives ont un $+A^4$ et les autres un $-A^4$.

Deux règles de signes sont énoncées, et il faut dire qu'elles ne s'accordent pas
tout à fait. La première : un côté au-dessous fait que « tout ce qui est affecté
sous lui » est aussi au-dessous. La seconde : lorsqu'on multiplie deux côtés de
nature différente, le produit est noté du signe $-$ et devient au-dessous, mais
« par équivalence » il doit être pris comme s'il était au-dessus. La première
règle, appliquée à la lettre, donne au carré d'un côté au-dessous le caractère
« au-dessous », ce que la seconde interdit ; et c'est de cette tension que naît
la difficulté du chapitre sur la quadratique qui subsiste par
elle-même.\footnote{« latus quod eiusmodi naturam sortitus est, efficere ut omnia
quæ sub ipso affecta sunt, sint etiam infra » (vue 4 L, premier point noté), et
« productum quod oritur notari quidem signo $-$, et fieri positivum infra, sed
per æquivalentiam (ut sic loquar) ita esse accipiendum ac si esset supra »
(même page). En notation moderne les deux règles se réduisent à une seule,
$(-1)\times(-1) = +1$, et il n'y a plus de tension ; le feuillet, qui tient deux
comptabilités distinctes — celle des signes écrits et celle du « dessus » et du
« dessous » —, en a deux à accorder.}

\section*{Ce que le traité énonce : les coefficients et les fonctions symétriques}

\pagerange{3}{9}
Au milieu du chapitre cinquième des cubiques, à la vue 9 L, le traité s'arrête
sur une remarque incidente et énonce, pour tous les degrés, la règle qui est
l'objet de tout son travail :

\begin{quote}
« in omnibus æquationibus alternatim affirmatis et negatis à principio usque ad
finem, coefficientem altioris gradus sub potestate fieri summam laterum,
coefficientem verò gradus immediate subsequentis summam rectangulorum vel
planorum sub iisdem lateribus […] similiter coefficientem gradus proximè
inferioris fieri summam omnium solidorum sub iisdem lateribus, et sic in
infinitum. »\footnote{« Dans toutes les équations alternativement affirmées et
niées du début jusqu'à la fin, le coefficient du degré le plus élevé sous la
puissance est la somme des côtés, le coefficient du degré immédiatement suivant
la somme des rectangles ou plans sous ces mêmes côtés, quels que soient les
côtés vrais qui ont été multipliés ; semblablement le coefficient du degré
immédiatement inférieur est la somme de tous les solides sous ces mêmes côtés,
et ainsi à l'infini. » Vue 9 L.}
\end{quote}

C'est, dans le langage des figures, ce que nous écrivons aujourd'hui ainsi. Si
$B_1, \dots, B_n$ sont les racines et
\[ P(A) = (A - B_1)(A - B_2)\cdots(A - B_n), \]
alors
\[ P(A) = A^n - e_1 A^{n-1} + e_2 A^{n-2} - \cdots + (-1)^n e_n, \]
où $e_k$ est la somme de tous les produits de $k$ racines distinctes : $e_1$ la
somme des racines, $e_2$ la somme de leurs produits deux à deux, $e_3$ la somme
de leurs produits trois à trois. Les $e_k$ s'appellent les \emph{fonctions
symétriques élémentaires} des racines, et la suite d'égalités s'appelle les
\emph{relations entre les coefficients et les racines}.\footnote{Les deux noms
sont postérieurs au feuillet et ne sont attachés ici qu'à cause de la
coïncidence exacte des énoncés : « somme des côtés » est $e_1$, « somme des
rectangles sous les mêmes côtés, pris deux à deux » est $e_2$, « somme de tous
les solides » est $e_3$. Le feuillet dit « rectangles ou plans » et « solides »
là où nous disons produits, parce que l'homogénéité lui impose de nommer la
dimension. Il faut ajouter deux réserves. La première est que le traité
\emph{vérifie} la règle degré par degré, en développant les produits, jusqu'au
degré quatre, et qu'il l'\emph{affirme} au-delà (« et sic in infinitum ») sans
la démontrer : il n'a ni récurrence, ni notation pour un degré indéterminé. La
seconde est plus lourde et fait l'objet du reste de cette section.}

Le traité a donc l'énoncé. Ce qu'il n'a pas, c'est le droit de s'en servir sans
exception, et il le sait. À la même page, il vient de construire une cubique en
multipliant la quadratique $G + A^2 = 0$ — celle qui « subsiste par elle-même »,
et dont on verra plus bas qu'elle n'a pas de racine réelle — par la latérale
$C - A = 0$. Le produit est
\[ GC - GA + CA^2 - A^3 = 0, \qquad \text{c'est-à-dire} \qquad A^3 - CA^2 + GA - GC = 0, \]
et il se range, par ses signes, dans la forme
$T\,\mathrm{Sol.} - z\,\mathrm{pl.}\,A + xA^2 - A^3 = 0$ qui est exactement celle
du premier chapitre, où les trois côtés sont au-dessus. Le traité voit aussitôt
que la forme est la même et l'interprétation différente, et il écrit la phrase
qui limite sa propre règle :

\begin{quote}
« $x$ enim ut videre est ex eius origine non est summa laterum, nec
$z\,\mathrm{pl.}$ aggregatum rectangulorum et cætera, ac proinde hæc constitutio
est exceptio regulæ generalis. »\footnote{« Car $x$, comme on le voit par son
origine, n'est pas la somme des côtés, ni $z\,\mathrm{pl.}$ l'agrégat des
rectangles, etc. ; et par conséquent cette constitution est une exception à la
règle générale. » Vue 9 L.}
\end{quote}

L'observation est exacte à l'intérieur de la doctrine : il n'y a ici qu'un seul
côté vrai, $C$, et $x$ vaut $C$ ; il n'y a aucun rectangle sous des côtés vrais,
et $z$ vaut $G$. Mais l'exception n'existe pas. Les trois racines de
$A^3 - CA^2 + GA - GC = 0$ sont $C$, $i\sqrt{G}$ et $-i\sqrt{G}$, et l'on vérifie
sur-le-champ :
\[ e_1 = C + i\sqrt{G} - i\sqrt{G} = C = x, \]
\[ e_2 = C\,i\sqrt{G} - C\,i\sqrt{G} + \bigl(i\sqrt{G}\bigr)\bigl(-i\sqrt{G}\bigr) = G = z, \]
\[ e_3 = C \cdot i\sqrt{G} \cdot \bigl(-i\sqrt{G}\bigr) = CG = T. \]
La règle générale vaut donc sans exception, et l'exception que le feuillet croit
trouver est le prix de son refus de compter les côtés fictifs comme des
racines.\footnote{Ce calcul est celui de cette lecture ; le feuillet le nie en
propres termes. C'est le seul endroit de tout le volume où l'on puisse montrer,
sur l'exemple même du feuillet, ce que coûte la doctrine et ce qu'elle gagnerait
à admettre les racines complexes. On verra à la dernière section du traité que
l'auteur en fournit lui-même une seconde démonstration, chiffrée, et tout aussi
concluante contre lui.}

Ce n'est pas une petite affaire, parce que toute la méthode repose sur la
possibilité de reconnaître une équation à ses signes. Or deux équations de
formes canoniques identiques peuvent avoir des interprétations tout à fait
différentes, et le traité le répète chapitre après chapitre — « quamvis omnino
similis sit formulæ quæ in primo capite, non tamen similiter enuntiabitur »,
« maximè tamen inter se differunt, ut patet ex earum diversa origine ». Il n'a,
pour les distinguer, aucun critère tiré de l'équation elle-même : il faut
connaître l'origine. La question est posée en trois lignes, sur un béquet
épinglé à un feuillet blanc, et l'on y reviendra : c'est la seule chose que la
doctrine ne sait pas faire, et la réponse moderne est le discriminant.

\section*{Les équations quadratiques}

\pagerange{4}{5}
Le traité commence par le degré deux, « parce que de toutes c'est la moins
composée », la latérale ne présentant aucune difficulté. Deux racines, donc
trois cas : toutes deux au-dessus, toutes deux au-dessous, ou l'une au-dessus et
l'autre au-dessous ; plus un quatrième chapitre pour la quadratique qui ne vient
d'aucun produit. Le traité remarque, en passant, que cette énumération est
légitime « puisque le nombre des cas est fini », ce qui est la justification
correcte d'une preuve par disjonction exhaustive.

Les produits sont écrits, et les formes canoniques qui en sortent sont, en
écriture moderne et monique :

\begin{itemize}
  \item \textbf{Caput 1}, deux racines positives $B$ et $C$ :
    $(A-B)(A-C) = A^2 - xA + z = 0$ avec $x = B + C$ et $z = BC$.
  \item \textbf{Caput 2}, deux racines négatives : $A^2 + xA + z = 0$, mêmes
    $x$ et $z$.
  \item \textbf{Caput 3}, une racine positive $B$ et une négative $-C$ :
    $(A-B)(A+C) = A^2 + (C - B)A - BC = 0$. Le traité en fait trois
    propositions selon que $B > C$, $B = C$ ou $B < C$, parce que dans son
    écriture le coefficient du terme en $A$ est \emph{la différence} des côtés,
    affectée du signe du plus grand ; le cas $B = C$ donne
    $A^2 - BC = 0$, où « l'affection sous le côté s'évanouit ».
  \item \textbf{Caput 4}, la quadratique qui subsiste par elle-même :
    $A^2 + z = 0$.
\end{itemize}

Le traité note en outre, et c'est juste, que la forme canonique ne dépend pas de
la complication apparente des coefficients : dans
\[ A^2 - \frac{KN}{M}A + \frac{HL}{G}A + \frac{F^{\mathrm{sol}}}{G} = 0, \]
les deux termes en $A$ « s'entendent comme une seule ligne, qu'on appelle $x$ »,
et l'on est ramené au premier chapitre. C'est la remarque qu'une forme canonique
est une forme, non une apparence.

\subsection*{La condition du discriminant}

Vient alors la doctrine des « déterminations », qui est la meilleure partie du
traité. Deux choses distinctes y sont appelées du même nom et il faut les
séparer, car le traité les sépare.

La première est le cas des racines égales. « Une équation est dite sujette à
détermination lorsqu'elle est explicable de deux ou plusieurs côtés égaux » ; la
détermination est \emph{majeure} si tous les côtés sont égaux, \emph{mineure} si
quelques-uns seulement le sont. Pour la quadratique du premier chapitre, cela
arrive « lorsque $x$ est coupé en deux parties égales et que chaque moitié est
$A$ » : les deux racines valent $x/2$, « à raison de quoi le côté est double et
la solution double aussi ». En termes modernes, la racine double est
$A = x/2$, valeur en laquelle $P$ et $P'$ s'annulent ensemble.

La seconde est la \emph{limitation}, et c'est la découverte : il y a une
inégalité en dehors de laquelle « les côtés seront feints, et par conséquent la
question impossible ». Le traité l'énonce et la démontre :

\begin{quote}
« $z\,\mathrm{pl.}$ non debet esse maius quadrante quadrati $x$, cum enim $x$ sit
summa laterum, rectangulum maius eo non potest quam si utraque latera sint
æqualia inter se. »\footnote{« $z\,\mathrm{pl.}$ ne doit pas être plus grand que
le quart du carré de $x$, car $x$ étant la somme des côtés, le rectangle ne peut
être plus grand que si les deux côtés sont égaux entre eux. » Vue 5 L. Le mot
rendu par « quadrante », le quart, est donné comme lecture incertaine par la
transcription ; les mathématiques le confirment sans réserve, la borne étant
$x^2/4$ et non une autre, et la phrase qui suit calculant explicitement « la
quatrième partie du carré de $x$ ».}
\end{quote}

C'est-à-dire : $z \le x^2/4$, ou encore
\[ x^2 - 4z \ge 0. \]
La quantité $x^2 - 4z$ est ce que nous appelons le \emph{discriminant} de
$A^2 - xA + z$, et sa positivité est exactement la condition pour que les deux
racines soient réelles.\footnote{Le nom est postérieur au feuillet ; l'énoncé
coïncide sans réserve, et la démonstration du feuillet est correcte : elle
consiste à voir que le produit de deux nombres de somme donnée est maximal
quand ils sont égaux, ce que nous démontrons aujourd'hui par
$BC = \tfrac14\bigl((B+C)^2 - (B-C)^2\bigr)$. On retiendra que la borne est
atteinte, et atteinte précisément au cas de la racine double : la limitation et
la détermination majeure se rencontrent en un même point, ce que le feuillet ne
dit pas mais que ses deux énoncés impliquent.}

Le traité ajoute, et c'est exact, que rien de tel n'est à craindre au troisième
chapitre, « lorsque $x$ est la différence » : une somme et une différence étant
données, le cas est toujours possible. En effet $A^2 + (C-B)A - BC$ a pour
discriminant $(C-B)^2 + 4BC = (B+C)^2 \ge 0$, toujours positif.

\subsection*{La quadratique qui subsiste par elle-même, et le nombre de ses racines}

Le quatrième chapitre est le point faible de tout l'ouvrage et il vaut la peine
de dire exactement pourquoi. Le traité écrit $z\,\mathrm{pl.} + A^2 = 0$ et
conclut :

\begin{quote}
« In quâ æquatione evidens est unicum esse latus positivum infrà potentia æquale
$z\,\mathrm{pl.}$, nam per antithesim erit $z\,\mathrm{pl.} = -A^2$. »\footnote{« Dans
laquelle équation il est évident qu'il y a un unique côté positif au-dessous,
égal en puissance à $z\,\mathrm{pl.}$ ; car par transposition on aura
$z\,\mathrm{pl.} = -A^2$. » Vue 5 L. « Égal en puissance à » veut dire : dont le
carré est.}
\end{quote}

Ce qui est vrai est ceci. L'équation $A^2 + z = 0$, avec $z$ positif, n'a
\emph{aucune} racine réelle : son discriminant vaut $-4z < 0$. Ses deux racines
sont $+i\sqrt{z}$ et $-i\sqrt{z}$, deux complexes conjuguées, dont le module
commun est $\sqrt{z}$. Le « côté unique au-dessous, égal en puissance à
$z\,\mathrm{pl.}$ », c'est-à-dire $\sqrt{z}$ affecté du caractère
« au-dessous », n'est pas une racine de l'équation : c'est le module de ses deux
racines. Le traité tient donc, sous le nom d'un côté, la seule quantité réelle
que la paire conjuguée fournisse, et il la fait fonctionner correctement partout
où il s'en sert — dans les chapitres 6, 7, 8 et 9 des quarto-quadratiques, où le
facteur $A^2 + G$ est toujours désigné par « le côté égal en puissance à
$G\,\mathrm{pl.}$ ». Cette lecture propose donc de traduire uniformément cette
expression par : \emph{le facteur quadratique $A^2 + G$, dont les deux racines
sont $\pm i\sqrt{G}$}.\footnote{La traduction est de cette lecture et va au-delà
de ce que le feuillet dit ; elle est proposée parce qu'elle rend cohérent tout
l'usage que le traité fait de cette quantité, et parce qu'elle est
mathématiquement exacte là où la lettre du feuillet ne l'est pas. Il faut
ajouter que le feuillet se contredit sur ce point : son dernier chapitre pose
comme critère qu'une quadratique $D\,\mathrm{pl.} + FA + A^2 = 0$ a ses côtés
feints dès que $D\,\mathrm{pl.} > \tfrac14 F^2$ ; appliqué à
$z\,\mathrm{pl.} + A^2 = 0$, où $F = 0$, ce critère donne aussitôt $z > 0$, donc
des côtés feints — et non l'unique côté au-dessous que le chapitre 4 lui
attribue. Les deux passages ne peuvent pas être vrais ensemble, et c'est le
dernier chapitre qui a raison.}

Une conséquence générale mérite d'être tirée, parce qu'elle empêche de prêter au
traité plus qu'il ne donne. Le traité ne dit \emph{pas}, et ne pouvait pas dire,
qu'une équation de degré $n$ a $n$ racines. Sa classification laisse au
contraire, à chaque degré, une classe résiduelle — les équations « qui
subsistent par elles-mêmes » — qui est précisément celle des équations ayant
moins de racines réelles que leur degré, et pour lesquelles il n'exhibe rien.
Une cubique à une seule racine réelle ne peut, dans son système, ni venir de
trois latérales, ni d'une quadratique à deux racines et d'une latérale : elle
subsiste par elle-même, et le chapitre 6 des cubiques n'en donne que les six
formes, sans racines. Le théorème selon lequel tout polynôme de degré $n$ à
coefficients complexes a exactement $n$ racines comptées avec leur multiplicité
— ce que nous appelons le théorème fondamental de l'algèbre — ne se trouve donc
pas sur ces feuillets, et cette lecture ne le leur attribue
pas.\footnote{L'énoncé et sa démonstration sont très postérieurs au feuillet. Ce
que le feuillet a, et qui est réel, c'est l'affirmation qu'une équation contient
toutes les racines dont elle a été composée, et l'inventaire complet, aux degrés
deux, trois et quatre, des manières de la composer avec des racines réelles.}

\section*{Les équations cubiques : les six chapitres}

\pagerange{6}{8}
Une cubique est explicable de trois côtés, d'où quatre chapitres selon les
signes ; et comme elle peut aussi naître d'une quadratique et d'une latérale, ou
subsister par elle-même, il y a deux chapitres de plus. Six en tout. Voici,
en écriture moderne monique, le catalogue entier, avec $x$, $z$, $T$ pour les
trois fonctions symétriques des côtés vrais quand elles ont ce sens.

\begin{itemize}
  \item \textbf{Caput 1}, trois racines positives :
    $A^3 - xA^2 + zA - T = 0$, avec $x = B+C+D$, $z = BC+BD+CD$, $T = BCD$.
  \item \textbf{Caput 2}, deux positives $B$, $C$ et une négative $-D$ :
    $(A-B)(A-C)(A+D) = 0$. Cinq propositions, selon la comparaison de $D$ à
    $B+C$ puis celle de $(B+C)D$ à $BC$ ; les formes sont
    $A^3 \pm xA^2 \pm zA - T = 0$ dans toutes les combinaisons de signes que
    ces comparaisons autorisent, le coefficient de $A^2$ étant $D - (B+C)$ et
    celui de $A$ étant $BC - (B+C)D$.
  \item \textbf{Caput 3}, deux négatives et une positive : le chapitre est
    l'inverse du précédent, et le traité le dit en une phrase — « toutes et
    chacune sont inverses, et le chapitre entier est l'inverse du chapitre
    entier ; les côtés qui là étaient au-dessus sont ici au-dessous ». C'est le
    changement de variable $A \mapsto -A$, correctement identifié et
    correctement exploité pour n'avoir rien à redémontrer.
  \item \textbf{Caput 4}, trois négatives : $A^3 + xA^2 + zA + T = 0$. Le
    traité remarque que toutes les affections y portent le signe $+$, « ce qui
    n'arrive jamais directement et par soi » — c'est-à-dire qu'une équation à
    tous ses coefficients positifs n'a aucune racine positive, ce qui est vrai
    et immédiat.
  \item \textbf{Caput 5}, née d'une quadratique et d'une latérale. Le traité
    écarte d'abord le cas où la quadratique a deux racines, qui retombe dans
    les chapitres précédents, et ne garde que la quadratique sans racine
    réelle : $(A^2+G)(A \mp C) = 0$, deux propositions.
  \item \textbf{Caput 6}, les cubiques qui subsistent par elles-mêmes : six
    formes, sans racines exhibées,
    \[ A^3 + xA^2 + T = 0, \quad A^3 + zA + T = 0, \quad A^3 + T = 0, \]
    \[ A^3 + zA - T = 0, \quad A^3 + xA^2 - T = 0, \quad A^3 - T = 0, \]
    les trois premières ayant « un unique côté au-dessous », les trois dernières
    « un unique côté au-dessus ».\footnote{Le feuillet les écrit sous la forme
    $T\,\mathrm{Sol.} + xA^2 + A^3 = 0$ et le reste, tous les termes au même
    membre. Les six formes sont exactement les cubiques réelles ayant une seule
    racine réelle et une paire complexe conjuguée, écrites selon le signe de
    cette racine et selon celui des coefficients qui ne s'annulent pas. Le
    feuillet ne dit pas que les deux autres racines sont feintes : il se borne à
    ne compter qu'un côté. La lecture le dit, parce que c'est ce qui est vrai.}
\end{itemize}

Lu ainsi, le classement du traité est une classification complète des cubiques
réeles moniques par le nombre et le signe de leurs racines réelles, jointe à la
lecture de chaque coefficient comme fonction symétrique de ces racines lorsque
toutes sont réelles. C'est, à un vocabulaire près, ce qu'un cours moderne fait
encore en préalable à l'étude du discriminant.

Trois petites choses sont à relever dans l'exécution, parce qu'elles sont dans
le feuillet et qu'un lecteur les rencontrera. À la vue 6 R, la forme canonique
du premier chapitre est écrite « $+\dot{x}A$ » sans l'exposant 2, que porte la
copie faite en marge par la seconde plume ; le terme est $xA^2$.\footnote{Vue
6 R. La transcription laisse la ligne telle quelle et le signale ; la marge
donne la forme juste.} Dans plusieurs produits, les termes de même degré sont
réunis par une accolade au lieu d'être écrits à la file, ce qui est une
commodité d'écriture et non une notation. Et le traité écrit une fois le côté
$a$ en minuscule là où il écrit $A$ partout ailleurs, sans conséquence.

\section*{Les déterminations des cubiques, et les limites du possible}

\pagerange{8}{11}
Cette partie est la plus originale du traité. Elle répond à deux questions
distinctes : comment trouver les racines quand deux d'entre elles sont égales,
et à quelle condition sur les coefficients les racines sont réelles.

\subsection*{La racine double, obtenue par éliminations successives}

Le traité prend la cubique du premier chapitre, $A^3 - xA^2 + zA - T = 0$, à
trois racines positives, et suppose que deux d'entre elles, $B$ et $C$, sont
égales. Il pose $B = C = A$, d'où le troisième côté $D = x - 2A$, et il traduit
les deux autres fonctions symétriques :
\[ z = A^2 + 2A(x - 2A) = 2xA - 3A^2, \qquad T = A^2(x - 2A) = xA^2 - 2A^3. \]
Ces deux égalités sont sur le feuillet et elles sont justes. Il en tire deux
équations en $A$ toutes deux nulles,
\[ \tfrac13 z - \tfrac23 xA + A^2 = 0, \qquad \tfrac12 T - \tfrac12 xA^2 + A^3 = 0, \]
multiplie la première par $A$ pour les rendre « homogènes », les compare,
retranche $A^3$ des deux membres, et obtient
\[ \tfrac12 T - \tfrac13 zA + \tfrac16 xA^2 = 0, \qquad \text{puis} \qquad \frac{3T}{x} - \frac{2zA}{x} + A^2 = 0 \]
en multipliant par $6$ et en divisant par $x$.\footnote{Le feuillet dit « quorum
omnium si sumatur sexta pars et productum applicetur ad $x$ » — « si de tout
cela on prend la sixième partie et qu'on applique le produit à $x$ ». L'opération
réellement faite est une multiplication par six suivie d'une division par $x$ ;
« prendre la sixième partie » est une ellipse pour « faire que le coefficient
$\tfrac16$ devienne l'unité ». Le résultat écrit est juste.} Une seconde
comparaison, avec la première équation, fait disparaître $A^2$ et donne enfin
\[ \frac{\dfrac{3T}{x} - \dfrac{z}{3}}{\dfrac{2z}{x} - \dfrac{2x}{3}} = A. \]

C'est la formule du feuillet, telle qu'il l'écrit. Mise sur dénominateur commun,
elle devient
\[ A = \frac{xz - 9T}{2\,(x^2 - 3z)}, \]
et cette expression est exacte : la racine double d'une cubique monique
$A^3 - xA^2 + zA - T$ qui en possède une est bien donnée par cette fraction
rationnelle des coefficients.\footnote{Vérification de cette lecture, sur les
racines $2, 2, 5$ : $x = 9$, $z = 24$, $T = 20$, d'où
$(xz - 9T)/2(x^2-3z) = (216-180)/2(81-72) = 36/18 = 2$. Et sur $3, 3, 1$ :
$x = 7$, $z = 15$, $T = 9$, d'où $(105-81)/2(49-45) = 24/8 = 3$. Les deux
formes, celle du feuillet et celle-ci, sont identiques et non seulement
équivalentes.}

Ce que le traité a fait là porte aujourd'hui un nom. Ses « comparaisons
successives » d'équations toutes égales à zéro, où l'on retranche à chaque pas
la plus haute puissance commune, sont exactement l'algorithme d'Euclide appliqué
au polynôme $P$ et à sa dérivée $P'$ ; et la fraction obtenue est la racine du
reste linéaire de cet algorithme. On peut le vérifier en deux lignes : avec
$P = A^3 - xA^2 + zA - T$ et $P' = 3A^2 - 2xA + z$, la division de $P$ par $P'$
laisse le reste $-\tfrac13 xA^2 + \tfrac23 zA - T$, et la combinaison suivante
laisse $2(x^2-3z)A - (xz - 9T)$, qui s'annule pour la valeur ci-dessus. Une
racine double de $P$ étant une racine commune à $P$ et $P'$, elle est racine de
leur plus grand commun diviseur, et c'est celui-ci que l'algorithme
calcule.\footnote{Le nom d'algorithme d'Euclide appliqué à $P$ et $P'$, et la
notion de dérivée, sont postérieurs au feuillet, qui n'a ni l'un ni l'autre :
il obtient ses deux équations de départ non par dérivation mais en traduisant la
condition $B = C$ dans les fonctions symétriques. La coïncidence est néanmoins
complète, et elle porte aussi sur la généralité : le traité affirme que « dans
toutes les constitutions, cubiques, quarto-quadratiques ou de degré plus élevé,
le côté égal se découvrira par le même artifice, sinon qu'il faudra établir
plusieurs équations pour parvenir enfin à la dernière où l'on aura $A$ unique ».
Cette affirmation est vraie, et c'est la description exacte de ce que fait
l'algorithme sur un polynôme de degré quelconque.}

\subsection*{Les limitations : ce que le traité établit, et ce qu'il croit établir}

Pour les deux chapitres où $x$ est la somme des trois côtés — le premier et le
quatrième —, le traité donne deux bornes et une raison :

\begin{quote}
« eo casu quæstio erit possibilis si $Z\,\mathrm{pl.}$ non sit maius tertia parte
$x$ quadrati, et similiter $T\,\mathrm{pl.}$ non sit maius tertiâ parte cubi
eiusdem $x$, si enim linea in tres partes dividatur maius solidum sub his
tribus portionibus contentum, erit cum dictæ tres partes erunt inter se
æquales. »\footnote{« Dans ce cas la question sera possible si
$z\,\mathrm{pl.}$ n'est pas plus grand que la troisième partie du carré de $x$,
et semblablement si $T$ n'est pas plus grand que la troisième partie du cube du
même $x$ ; car si une ligne est divisée en trois parties, le plus grand solide
contenu sous ces trois portions sera lorsque lesdites trois parties seront
égales entre elles. » Vue 8 L. Le feuillet écrit ici « $T\,\mathrm{pl.}$ », un
plan, là où il s'agit du solide $T\,\mathrm{Sol.}$ : le mot de dimension est
fautif, et la loi d'homogénéité suffit à le voir.}
\end{quote}

La première borne est juste. Si $B + C + D = x$ avec $B$, $C$, $D$ positifs,
alors $z = BC + BD + CD$ est maximal quand les trois sont égaux à $x/3$, et vaut
alors $3(x/3)^2$ :
\[ z \le \frac{x^2}{3}. \]

La seconde, telle qu'elle est écrite, est fausse. Le raisonnement que le
feuillet donne — le solide sous trois parties d'une ligne est maximal quand les
parties sont égales — est correct, et il conduit non pas au tiers du cube mais
au cube du tiers :
\[ T \le \left(\frac{x}{3}\right)^3 = \frac{x^3}{27}. \]
La borne vraie est donc vingt-sept fois plus petite que celle que le feuillet
énonce. Il est probable que « la troisième partie du cube » soit une ellipse
pour « le cube de la troisième partie », d'autant que la phrase parallèle sur
$z$ — « la troisième partie du carré » — est, elle, littéralement juste ; mais
telle qu'elle est écrite la borne ne l'est pas, et cette lecture écrit
$x^3/27$.\footnote{Départ de cette lecture par rapport au feuillet : la borne
énoncée est $x^3/3$, la borne vraie est $x^3/27$, et la démonstration du
feuillet donne la seconde. L'erreur est sans conséquence sur la suite, parce que
le traité ne se sert jamais de cette borne pour calculer : il s'en sert pour
dire qu'il y a une borne.}

Reste le plus important, que le feuillet n'a pas vu. Ces deux inégalités sont
\emph{nécessaires} pour que les trois racines soient réelles et positives ;
elles ne sont pas \emph{suffisantes}, et le feuillet les prend pour telles en
écrivant « eo casu quæstio erit possibilis si… ». Un contre-exemple entier suffit
à le montrer :
\[ A^3 - 6A^2 + 10A - 8 = 0, \]
c'est-à-dire $x = 6$, $z = 10$, $T = 8$. Les deux limitations sont satisfaites,
$10 \le 36/3 = 12$ et $8 \le 216/27 = 8$ — et a fortiori la borne trop large du
feuillet, $8 \le 72$. Pourtant cette équation se factorise en
\[ (A - 4)\,(A^2 - 2A + 2) = 0 \]
et n'a qu'une seule racine réelle, $4$, les deux autres étant $1 + i$ et
$1 - i$. La question est donc « impossible » au sens du traité alors que ses
deux conditions sont remplies.\footnote{Calcul et contre-exemple de cette
lecture. Le critère exact est le signe du discriminant : pour
$A^3 - xA^2 + zA - T$,
\[ \Delta = 18xzT - 4x^3T + x^2z^2 - 4z^3 - 27T^2, \]
et les trois racines sont réelles et distinctes si et seulement si
$\Delta > 0$, deux d'entre elles complexes conjuguées si $\Delta < 0$, et il y a
une racine multiple si $\Delta = 0$. Sur l'exemple : $\Delta = -400$. Sur les
racines $1, 2, 3$ ($x=6$, $z=11$, $T=6$) : $\Delta = 4$, qui est bien
$\prod_{i<j}(B_i - B_j)^2$. Le mot et la quantité sont postérieurs au feuillet ;
on verra à la section suivante que les limitations que le traité obtient pour
les \emph{autres} chapitres sont, elles, exactement des conditions de
discriminant, et qu'il les obtient justes.}

\subsection*{Les déterminations chapitre par chapitre}

Le traité passe ensuite en revue les cinq propositions du second chapitre, et
son travail est propre. On ne relèvera ici que ce qui a une portée.

Pour la deuxième proposition, où le côté au-dessous égale la somme des deux
côtés au-dessus, la détermination se calcule ainsi : si $D = 2A$ et $B = C = A$,
alors $z = 3A^2$, donc $A = \sqrt{z/3}$, et $T = 2A^3$. Le traité le vérifie de
deux manières, par l'analyse directe puis au moyen d'un lemme emprunté, et les
deux donnent $A^2 = z/3$. Le calcul est juste : de $T = 2A^3$ et
$T = zA - A^3$ suit $3A^3 = zA$, donc $3A^2 = z$.

Pour la troisième proposition, trois cas selon que le côté au-dessous égale
chacun des deux autres, un seul, ou aucun. Dans le premier, la détermination est
majeure et $T = x^3$, $z = x^2$ ; dans le second, $z = T/x$ ; dans le troisième,
on revient à l'élimination générale.\footnote{Le premier cas est à lire avec
attention : les trois côtés sont alors égaux en valeur absolue, mais l'un est
au-dessous, en sorte que $x$ n'est plus la somme des trois mais leur différence,
et $T = x^3$, $z = x^2$ s'entendent de cette différence. Le feuillet dit
« $T\,\mathrm{Sol.}$ æquabitur $x^3$ sicut et $z\,\mathrm{pl.}$ quadrato ipsius
$x$ », sans plus, et cette lecture ne lui prête pas davantage.}

Pour la cinquième proposition, le traité démontre un lemme utile : si les deux
côtés au-dessus $B$, $C$ ont un rectangle $BC$ plus grand que $(B+C)D$, alors
chacun d'eux est plus grand que $D$. La démonstration est correcte —
$(B+C)C > BC > (B+C)D$ donne $C > D$, et de même pour $B$. La seconde plume a
écrit en marge, en regard du mot « maiorem » qu'elle a souligné dans le texte,
le mot « minorem » : c'est une correction, et elle est fausse. La démonstration
demande bien « plus grand », la ligne du feuillet est juste, et l'annotateur a
tort.\footnote{Vue 11 R. Le texte porte « similique ratione demonstrabitur etiam
rectam $AB$ maiorem esse ipsâ $D$ », la marge « minorem ». Par symétrie de $B$
et $C$ dans l'hypothèse $BC > (B+C)D$, la conclusion est la même pour l'un que
pour l'autre : les deux sont plus grands que $D$. C'est l'un des deux seuls
endroits du volume où l'annotateur se prononce sur la substance, et il se trompe.}

\section*{Les trois lemmes de maximum, et ce que « D. Rob. » veut dire}

\pagerange{9}{11}
Deux fois, pour établir ses limitations, le traité s'appuie sur des lemmes qu'il
ne démontre pas et dont il dit d'où ils viennent : « præmisso Lemmate quod a
D. Rob. demonstratum est » (vue 9 R), puis « quod demonstrabitur præmissis
duobus Lemmatibus etiam à D. Rob. demonstratis » (vue 10 R). En marge de la
première mention, une note en français, la seule note non latine de tout le
traité, renvoie le lecteur à la source : « 13 en son traité de mechanique des
plans inclinez ».

Il faut dire exactement ce que cela établit. Cela établit que le traité
\emph{cite} Roberval, comme l'auteur de trois lemmes qu'il emprunte et qu'il
signale comme empruntés, et qu'un lecteur — peut-être l'auteur lui-même — a noté
en marge où les trouver. Cela n'établit pas que le traité soit de Roberval :
c'est même le contraire qu'une telle mention indique, un auteur n'ayant pas
coutume de renvoyer à lui-même comme à un tiers nommé par l'initiale de sa
qualité. Cette lecture n'attribue donc le traité à personne — ni à Roberval, ni
à Viète, ni à un autre — et note seulement que son auteur avait lu Roberval sur
les plans inclinés.\footnote{Le traité est anonyme sur le feuillet ; aucune
souscription, aucun nom d'auteur, aucune date. Le catalogue range le volume sous
le nom de Roberval ; les feuillets ne le disent pas. Sur la note marginale : le
chiffre « 13 » qui la précède appartient à la série des petits chiffres courants
de la seconde plume et n'est pas un renvoi de page.}

Voici les trois lemmes, traduits, et ce qu'ils valent.

\subsection*{Premier lemme (vue 10 L) : le plus grand solide appliqué à un plan}

\begin{quote}
« De tous les parallélépipèdes appliqués selon des plans égaux, ou selon le même
plan parallélogramme, et déficients de figures parallélépipédiques semblables
entre elles et semblablement posées, le plus grand est celui qui est appliqué
aux deux tiers, étant double du défaut. »
\end{quote}

En termes modernes, il s'agit de maximiser
\[ f(A) = zA - A^3 \]
— un solide « appliqué au plan $z$ » et « déficient d'une figure cubique ». Le
maximum est atteint pour $f'(A) = z - 3A^2 = 0$, soit
\[ A = \sqrt{\frac{z}{3}}, \qquad f(A) = \frac{2z}{3}\sqrt{\frac{z}{3}}. \]
Tout ce que dit le lemme s'y retrouve exactement. Posons
$\ell = \sqrt{z/3}$, de sorte que $z = 3\ell^2$ : le maximum vaut $2\ell^3$, la
figure déficiente vaut $\ell^3$, et le solide est donc bien « double du
défaut » ; et comme $2\ell^3 = \tfrac23 z \cdot \ell$, il est « appliqué aux
deux tiers » du plan. Le corollaire que le traité en tire dit la même chose
autrement, et dissipe l'ellipse relevée plus haut : « tout plan est contenu sous
deux droites inégales, dont la moindre peut la troisième partie du même plan, et
la plus grande est triple de la moindre » — c'est-à-dire $z = \ell \cdot 3\ell$
avec $\ell^2 = z/3$. La « hauteur égale au tiers du plan » des pages précédentes
est donc, sans ambiguïté, la droite dont le carré est le tiers du plan, soit
$\sqrt{z/3}$ ; et le traité l'écrit lui-même en clair une page plus loin :
« manifestum est $A$ fieri $\sqrt{\tfrac13 z\,\mathrm{pl.}}$ ».\footnote{Vue 9 R.
La lecture de l'ellipse est donc donnée par le feuillet lui-même, et n'est pas
une conjecture. Le mot « potens » — « peut » — est le terme technique de l'école
pour « dont le carré égale » ; il vient du grec $\delta\acute{\upsilon}\nu\alpha\tau\alpha\iota$
des \emph{Éléments}.}

L'usage de ce lemme est la limitation de la deuxième proposition du deuxième
chapitre : $T$ ne doit pas être plus grand que le maximum, soit
\[ T \le \frac{2z}{3}\sqrt{\frac{z}{3}}, \qquad \text{c'est-à-dire} \qquad 27\,T^2 \le 4\,z^3. \]
Cette dernière inégalité est exactement la condition pour que la cubique réduite
$A^3 - zA + T = 0$ ait ses trois racines réelles : c'est la positivité de son
discriminant, $\Delta = 4z^3 - 27T^2$. Le traité l'atteint, et l'atteint juste,
par un problème de maximum au lieu d'un invariant
algébrique.\footnote{Coïncidence exacte, et le nom de discriminant est
postérieur au feuillet. C'est la meilleure observation mathématique de tout le
dossier : chaque « limitation » que le traité tire de ces lemmes est, en langage
moderne, la condition $\Delta \ge 0$ pour la cubique correspondante, obtenue en
cherchant le maximum de la fonction dont la limitation exprime que $T$ ne le
dépasse pas. Le procédé est général et il est correct.}

\subsection*{Deuxième et troisième lemmes (vue 10 R) : la ligne divisée, et le solide appliqué à une ligne}

\begin{quote}
« Si une ligne droite est divisée en deux parties inégales, le plus grand solide
qui puisse être fait du carré de l'une par la hauteur de l'autre est lorsque la
portion dont le carré est la base dudit solide est double de l'autre portion. »

« De tous les parallélépipèdes appliqués selon la même ligne droite, ou selon
des droites égales, et déficients de figures parallélépipédiques semblables
entre elles et semblablement posées, le plus grand est celui qui est appliqué à
la troisième partie, étant la moitié du défaut. »
\end{quote}

Ces deux lemmes sont le même énoncé sous deux habits, et il vaut de le dire,
puisque le traité s'en sert comme de deux et remarque ensuite que le second
n'est pas nécessaire. Maximiser $u^2 v$ sous la contrainte $u + v = x$, c'est
maximiser $xu^2 - u^3$ ; le maximum est en $u = \tfrac23 x$, $v = \tfrac13 x$, et
vaut
\[ \frac49 x^2 \cdot \frac{x}{3} = \frac{4}{27}\,x^3. \]
Le premier lemme est l'énoncé sous la forme « $u$ double de $v$ » ; le second
sous la forme « appliqué à la troisième partie », la hauteur du solide maximal
étant $x/3$, et « moitié du défaut », le défaut valant
$u^3 = \tfrac{8}{27}x^3$, exactement le double. Les deux énoncés sont justes et
équivalents.

L'usage est la limitation de la quatrième proposition du deuxième chapitre : on
y a $(B+C)D = BC$, et le traité démontre que
\[ T \le \frac{4}{27}\,x^3, \]
ce qu'il écrit en clair, « scilicet $\tfrac{4}{27}x^3$ ». Là encore, l'inégalité
est exactement la condition de réalité des racines de $A^3 - xA^2 + T = 0$, dont
le discriminant vaut $4x^3T - 27T^2 = T(4x^3 - 27T)$, positif précisément pour
$0 \le T \le \tfrac{4}{27}x^3$. Deuxième coïncidence exacte, par la même voie.

Le traité donne en outre une seconde démonstration de cette même limitation,
purement géométrique, qui n'a besoin que du premier lemme, et il la conduit en
deux cas sur la figure d'un segment $AC$ partagé en $AB$ et $BC$ avec un
segment $CD$ dans le prolongement. Cette lecture l'a refaite et elle est juste
dans les deux cas. Le premier cas, où $CD$ vaut le quart de $AC$ et où
$B$ est le milieu, donne l'égalité : $T = \tfrac14 AC^2 \cdot \tfrac14 AC =
\tfrac1{16}AC^3$, tandis que $x = \tfrac34 AC$ et $\tfrac4{27}x^3 =
\tfrac4{27}\cdot\tfrac{27}{64}AC^3 = \tfrac1{16}AC^3$ : les deux quantités sont
bien égales. Le second cas, où $CD$ est moindre que le quart, donne l'inégalité
stricte par comparaison de deux solides dont la base et la hauteur sont l'une et
l'autre plus grandes.\footnote{Une lettre de la démonstration est à corriger, et
les mathématiques du feuillet suffisent à le faire : la phrase « sit primo recta
$CD$ æqualis quartæ parti rectæ $AE$ » demande $AC$ et non $AE$, comme le montre
la suite immédiate de la phrase, « igitur et $AF$ est quarta pars quadrati rectæ
$AC$ », et comme le demande la construction, $AE$ n'étant défini que deux lignes
plus loin. La transcription donne d'ailleurs « $AE$ » comme lecture incertaine.
Cette lecture écrit $AC$.}

Un mot enfin sur la figure de la vue 10 L, parce que la note marginale qui
l'accompagne est le seul endroit où le copiste parle en son nom : « N\textsuperscript{a}.
quod linea $DB$ debet esse æqualis $CD$. sed defectu cartæ non potuit ulterius
produci » — « noter que la ligne $DB$ doit être égale à $CD$, mais faute de
papier elle n'a pu être poussée plus loin ». Le dessin est donc faux, et
sciemment ; la condition manquante est rétablie par la note.

\section*{Les équations quarto-quadratiques : dix chapitres}

\pagerange{12}{22}
Le degré quatre occupe plus de la moitié du traité, et il est construit sur le
même plan, avec dix chapitres annoncés d'avance selon la manière dont l'équation
naît :

\begin{itemize}
  \item de quatre latérales — cinq chapitres, selon le nombre de côtés au-dessus
    (quatre, trois, deux, un, aucun) ;
  \item d'une quadratique et de deux latérales ;
  \item de deux quadratiques ;
  \item d'une cubique à deux côtés et d'une latérale ;
  \item d'une cubique subsistant par elle-même et d'une latérale ;
  \item ou bien elle subsiste par elle-même, ou bien des côtés feints entrent
    dans sa composition.
\end{itemize}

Le premier chapitre donne la forme de référence. Quatre racines positives,
$(A-B)(A-C)(A-D)(A-F) = 0$, soit
\[ A^4 - xA^3 + zA^2 - TA + S = 0 \]
avec, et le traité l'énonce terme à terme, $x$ la somme des quatre côtés, $z$
l'agrégat des six rectangles deux à deux, $T$ la somme des quatre solides trois
à trois, $S$ le plano-plan sous les quatre. Ce sont $e_1$, $e_2$, $e_3$, $e_4$,
et le compte des termes — quatre, six, quatre, un — est juste.

Les chapitres 2 à 5 suivent, par changement de signes et par la remarque, faite
et refaite, qu'un chapitre est l'inverse d'un autre : le cinquième est l'inverse
du premier, le quatrième du deuxième. Dans chaque chapitre, les propositions
sont les configurations de signes que les comparaisons de grandeurs autorisent,
et le traité les obtient par un travail d'inégalités qui est le vrai contenu de
la partie. Pour n'en donner qu'un exemple entièrement vérifié : au deuxième
chapitre, trois côtés au-dessus $B$, $C$, $D$ et un au-dessous $-F$, si $F$ est
plus grand que $B + C + D$, alors le coefficient de $A^2$ est positif, parce que
\[ F(B+C+D) > (B+C+D)^2 > BC + BD + CD, \]
et le coefficient de $A$ est négatif, parce que $BCD$ est plus petit que chacun
de $BCF$, $BDF$, $CDF$. Les deux raisonnements sont sur le feuillet et ils sont
corrects.\footnote{Vue 12 R. Le feuillet dit « rectangulum sub $F$ et summâ
dictorum trium laterum maius erit quadrato eorundem, quod quidem quadratum patet
esse maius tribus rectangulis negatis » ; c'est bien l'inégalité
$(B+C+D)^2 > BC+BD+CD$, qui tient parce que le carré contient en plus
$B^2 + C^2 + D^2$ et un second exemplaire de chaque rectangle.}

Le troisième chapitre — deux côtés au-dessus, deux au-dessous — porte le seul
raisonnement long et difficile de cette partie, sur les vues 13 et 16, et cette
lecture l'a refait en entier. Il s'agit de montrer que si la somme des côtés
au-dessous surpasse celle des côtés au-dessus et si les rectangles affirmés
égalent les niés, alors les solides niés surpassent les affirmés. Notons $s$ la
somme des côtés au-dessus, $t$ celle des côtés au-dessous, $p$ le rectangle sous
les premiers, $q$ celui sous les seconds ; l'hypothèse est $p + q = st$ avec
$t > s$, et il faut montrer $qs > pt$. Le traité y parvient ainsi, et chaque pas
tient : de $p \le s^2/4$ et $p + q = st > s^2$ suit $q > \tfrac14 s^2 \ge p$,
donc $2q > p + q = st$, donc $4q > st$ ; d'où
$s/t = s^2/st > s^2/4q \ge p/q$, et enfin $qs > pt$. C'est une belle
démonstration, correcte, conduite entièrement en proportions.

Elle se termine pourtant sur une phrase qu'il faut redresser. Le traité ajoute :
« sed et facile demonstrabitur prius illud solidum maius esse septuplo
posterioris » — « mais on démontrera aussi facilement que ce premier solide est
plus grand que le septuple du second ». Ce n'est pas vrai, et ne peut pas
l'être. Sous les hypothèses posées, le rapport $qs/pt$ des deux solides est
toujours \emph{inférieur à 4}, et il est toujours supérieur à $2+\sqrt3 \approx
3{,}73$ ; il ne peut donc jamais atteindre 7. Un exemple entier le montre : les
deux côtés au-dessus étant $\tfrac12$ et $\tfrac12$, les deux au-dessous
$\tfrac32$ et $\tfrac52$, on a $s = 1$, $t = 4$, $p = \tfrac14$, $q = \tfrac{15}4$,
les rectangles affirmés $p + q = 4$ égalent les niés $st = 4$ comme il est
supposé, et le rapport des solides vaut $qs/pt = \tfrac{15}{4}$, soit
$3{,}75$.\footnote{Calcul et exemple de cette lecture. La borne supérieure se
voit en écrivant $qs/pt = s^2/p - s/t$, qui croît quand $p$ décroît et reste
inférieur à $4$ puisque $p \le s^2/4$ ; la borne inférieure est atteinte à la
limite $t = (2+\sqrt3)s$, $p = s^2/4$, $q = t^2/4$, seule configuration où les
deux contraintes $p \le s^2/4$ et $q \le t^2/4$ soient compatibles avec
$p + q = st$. Il faut ajouter que le mot « septuplo » est donné par la
transcription comme lecture incertaine : l'écart peut donc être de la
transcription et non de l'auteur. Ce qui est sûr, c'est qu'aucun multiple de 7
n'est possible, et qu'un multiple compris entre $2+\sqrt3$ et $4$ l'est. La
phrase est du reste une remarque incidente, annoncée comme facile et jamais
démontrée ; le raisonnement principal ne s'appuie pas sur elle.}

Deux béquets épinglés à l'angle d'un feuillet blanc, et photographiés l'un après
l'autre aux vues 14 et 15 tandis qu'on les soulevait, portent les deux seules
additions au texte. Le premier, intitulé « p. 21. », fournit une phrase qui
manquait, sur l'origine des deux rectangles affirmés du troisième chapitre : l'un
vient du produit des côtés au-dessus, l'autre du produit des côtés au-dessous.
Le second, intitulé « 19\textsuperscript{a} », est une question :

\begin{quote}
« quomodo cognoscetur utrum hæc æquatio $T\,\mathrm{Sol.} - z\,\mathrm{pl.}\,A +
xA^2 - A^3 = 0$ sit capitis quinti vel tertii. »\footnote{« Comment
connaîtra-t-on si cette équation est du cinquième chapitre ou du troisième. »
Vue 15. L'équation écrite est celle qui est commune au premier chapitre des
cubiques — trois racines positives — et à la première proposition du cinquième —
une quadratique sans racine réelle multipliée par une latérale ; le béquet nomme
le troisième chapitre là où les signes désignent le premier, et cette lecture ne
corrige pas la référence, faute de savoir si le chiffre est du feuillet ou de la
main qui l'a écrit.}
\end{quote}

C'est la question juste, et c'est celle à laquelle la doctrine ne peut pas
répondre : les deux constitutions ont les mêmes signes, et le traité n'a aucun
moyen de les distinguer sans connaître l'origine. La réponse moderne est le
signe du discriminant donné plus haut : positif, trois racines réelles, donc
trois côtés vrais ; négatif, une seule racine réelle et une paire complexe, donc
le cinquième chapitre. L'exemple $A^3 - 6A^2 + 10A - 8 = 0$ construit plus haut
est justement du cinquième chapitre — il est
$(A^2 - 2A + 2)(A - 4)$ — et son discriminant vaut $-400$.

\subsection*{Les chapitres 6 à 9 : les facteurs quadratiques sans racine réelle}

Les quatre derniers chapitres constructifs font entrer le facteur $A^2 + G$ dont
il a été question. Ils sont l'endroit où la doctrine travaille le mieux et où
elle est le plus près de basculer.

\begin{itemize}
  \item \textbf{Caput 6}, deux quadratiques sans racine réelle :
    $(A^2+G)(A^2+H) = A^4 + (G+H)A^2 + GH = 0$. Le traité écrit que $z$ est la
    somme des deux « carrés » et $S$ le plan sous eux, ce qui est juste, et que
    l'équation est explicable de deux côtés au-dessous, ce qui ne l'est pas :
    elle n'a aucune racine réelle, ses quatre racines étant
    $\pm i\sqrt{G}$ et $\pm i\sqrt{H}$.\footnote{Départ de cette lecture. Le
    feuillet dit « sunt duo latera infrà quorum quadrata simul sumpta exhibent
    $z\,\mathrm{pl.}$ » ; les deux quantités $\sqrt G$ et $\sqrt H$ dont les
    carrés donnent $z$ existent bien et sont les modules des deux paires
    conjuguées, mais ce ne sont pas des racines de l'équation.}
  \item \textbf{Caput 7}, une quadratique sans racine réelle et deux latérales :
    $(A^2+G)(A \mp B)(A \mp C) = 0$, onze propositions. Le traité énonce les
    coefficients avec exactitude : $x$ la somme (ou la différence) des deux
    côtés vrais, $z$ la somme du rectangle sous eux et du « carré » $G$, $T$ le
    produit de $G$ par la somme des deux côtés, $S$ le produit de $G$ par leur
    rectangle. Ces quatre égalités sont justes — et elles sont, exactement, les
    fonctions symétriques élémentaires des quatre racines $B$, $C$,
    $\pm i\sqrt G$, la paire conjuguée contribuant $0$ à $e_1$ et $G$ à
    $e_2$.\footnote{Vérification de cette lecture. C'est la troisième fois que
    le feuillet produit lui-même la preuve que sa règle générale vaut pour les
    côtés fictifs comme pour les autres, et la deuxième fois qu'il s'arrête au
    seuil de la conclusion.}
  \item \textbf{Caput 8}, une cubique à deux côtés et une latérale : vingt
    propositions, tirées des deux formes de cubique possibles,
    $(A^2+G)(A \mp F)$, multipliées chacune par une latérale de chaque signe.
  \item \textbf{Caput 9}, une cubique subsistant par elle-même et une latérale :
    vingt propositions, tirées des six formes du sixième chapitre des cubiques
    multipliées chacune par $A - B$ et par $A + C$.
\end{itemize}

Les douze produits du chapitre 9 ont été refaits un à un dans cette lecture, et
onze sont exacts. Le seul qui ne le soit pas est celui de $A^3 - R$ par $A - B$ :
le feuillet écrit $RB + RA - BA^3 + A^4 = 0$ là où le produit donne
\[ (R - A^3)(B - A) = RB - RA - BA^3 + A^4, \]
avec un signe $-$ devant $RA$. Le second cas, $(R - A^3)(C + A) = RC + RA - CA^3
- A^4$, est écrit juste avec un $+$.\footnote{Vue 22 L. La transcription relève
le signe et le laisse comme il est écrit. Il en va de même d'une omission plus
bénigne : à la vue 21 L, l'équation $R\,\mathrm{sol.} \; G\,\mathrm{pl.}\,A +
A^3 = 0$ est écrite sans aucun signe entre les deux premiers termes, le trait
d'abréviation de « sol. » se prolongeant à droite ; le produit calculé deux
lignes plus bas suppose un $+$, et cette lecture écrit $A^3 + GA + R = 0$.}

Le chapitre 9 laisse aussi un trou. Sa septième proposition est énoncée — « cum
latus supra minus est latere infra » — et aucune formule n'est écrite au-dessous,
alors que la phrase suivante parle « de ces trois formules ». La formule
manquante se déduit du feuillet lui-même, sans rien emprunter : les propositions
5, 6 et 7 sortent du même produit et ne diffèrent que par la comparaison des
deux côtés, laquelle gouverne le coefficient $x$, leur différence ; la
proposition 5 a $+xA^3$, la proposition 6, où les côtés sont égaux, n'a plus de
terme en $A^3$, et la proposition 7 doit donc avoir
\[ S - TA + zA^2 - xA^3 - A^4 = 0. \]
C'est l'inférence de cette lecture, et elle ne repose que sur la structure des
deux propositions voisines.\footnote{Aucune édition imprimée n'a été consultée
pour la combler, et cette lecture ne prétend pas restituer ce que le feuillet
portait : elle dit ce que le produit et la série exigent. La même série fournit
la vérification : la proposition 6, où le terme en $A^3$ disparaît parce que la
différence des côtés est nulle, confirme que c'est bien $x$ qui porte cette
comparaison.}

Deux remarques du chapitre 9 méritent d'être relevées parce qu'elles sont justes
et fines. Le traité observe que dans ces formules $x$ n'est plus la somme de
plusieurs côtés mais un seul côté, celui de la latérale, et qu'en conséquence il
ne s'évanouit pas même quand les deux côtés sont égaux :

\begin{quote}
« Ideo autem in hac formula etiam si latera sint æqualia, et $x$ non sit eorum
summa, non tamen idem $x$ evanescit, quia ut patet ex origine istius est æquale
lateri infra æquationis lateralis. »\footnote{« Or dans cette formule, même si
les côtés sont égaux et que $x$ ne soit pas leur somme, ce même $x$ ne
s'évanouit pourtant pas, parce que, comme il ressort de son origine, il est égal
au côté au-dessous de l'équation latérale. » Vue 22 L.}
\end{quote}

L'observation est exacte, et la raison moderne en est transparente : les trois
racines de la cubique $A^3 + GA \pm R = 0$, qui n'a pas de terme en $A^2$, ont
une somme nulle ; la somme des quatre racines du produit est donc celle de la
seule racine de la latérale. Le feuillet voit l'effet et en donne la bonne
raison — « son origine » — sans avoir la phrase sur la somme nulle des trois
racines, qui demanderait de compter les fictives.

Un dernier point de tenue, qui vaut d'être dit pour qu'un lecteur ne s'y égare
pas. Cette lecture a vérifié tous les produits, c'est-à-dire toutes les
multiplications que le traité effectue. Elle n'a pas revérifié un à un les
quatre-vingt-sept motifs de signes des formules interprétées, et il faut dire
que tous ne suivent pas de leur produit tels qu'ils sont écrits : au huitième
chapitre, les propositions 2 à 10 portent un $+A^4$ alors que le produit d'où
elles sortent se termine par $-A^4$. En écrivant les équations moniques, comme
le fait cette lecture, la question disparaît, le signe du terme dominant
n'étant plus qu'une convention d'écriture.\footnote{La transcription relève de
son côté deux désaccords entre la ligne et la forme recopiée en marge par la
seconde plume : à la vue 12 R, proposition 2, la ligne porte $+A^4$ et la marge
$-A^4$, et c'est la marge qui s'accorde avec la constitution ; à la vue 20 L,
proposition 5, la ligne porte $+z\,\mathrm{pl.}\,A^2$ et la marge
$-z\,\mathrm{pl.}\,A^2$. Les deux sont donnés tels quels de part et d'autre dans
la transcription.}

\section*{Le dernier chapitre : les côtés feints, et l'exemple chiffré}

\pagerange{23}{23}
Le dixième chapitre, « et dernier », traite des quarto-quadratiques qui
subsistent par elles-mêmes « ou dans la composition desquelles des côtés feints
se mêlent ». C'est le meilleur chapitre du traité, et il porte sa propre
réfutation.

Le traité y pose d'abord le critère, et le pose juste. Soit
$A^2 - CA + B = 0$ : elle est explicable de deux côtés au-dessus, dont la somme
est $C$, « et elle est possible dans le cas où $B$ ne sera pas plus grand que la
quatrième partie de $C^2$ ». Si, sous la même forme, on construit une équation
où $B$ est plus grand que ce quart, « il est certain que les deux côtés dont
cette dernière équation naît sont feints ». C'est-à-dire : les racines de
$A^2 - CA + B$ sont réelles si et seulement si $C^2 - 4B \ge 0$. Le traité tient
donc parfaitement le discriminant de la quadratique et sait s'en servir pour
reconnaître une paire complexe.

Il construit alors, exprès, le produit d'une quadratique à racines réelles par
une quadratique à racines feintes, et il le fait en nombres. Ses deux
quadratiques sont
\[ A^2 - 6A + 5 = 0 \qquad \text{et} \qquad A^2 + 3A + 20 = 0, \]
la première de racines $1$ et $5$, la seconde sans racine réelle puisque
$20 > 9/4$, ce que le traité relève : « patet in ultima æquatione latera esse
ficta, ideoque eam impossibilem ». Le produit, qu'il écrit dans la notation
cossique, est
\[ A^4 - 3A^3 + 7A^2 - 105A + 100 = 0, \]
et cette lecture l'a vérifié : le développement est exact, terme à terme. Le
traité vérifie ensuite que $5$ est bien racine, en égalant les affirmés aux
niés : $100 + 7\cdot 25 + 625 = 900$ d'un côté, $105 \cdot 5 + 3 \cdot 125 = 900$
de l'autre. Les deux colonnes du feuillet donnent $900$ et $900$, et le calcul
est juste. Il ajoute, sans le refaire, que l'égalité vaut de même pour $1$ ; elle
vaut en effet.

Mais le traité conclut, une page plus haut, que dans une telle formule « $x$ est
la somme des deux côtés dont l'équation est explicable ». Sur son propre exemple,
cela est faux, et de façon vérifiable en une ligne. Les deux côtés vrais sont
$1$ et $5$ ; leur somme est $6$. Le coefficient du terme en $A^3$ est $3$, non
$6$. Ce qui vaut $3$, c'est la somme des \emph{quatre} racines : $1$, $5$, et les
deux racines de $A^2 + 3A + 20$, qui sont
\[ \frac{-3 + i\sqrt{71}}{2} \qquad \text{et} \qquad \frac{-3 - i\sqrt{71}}{2}, \]
dont la somme est $-3$. Et $1 + 5 - 3 = 3$. De même le terme constant, $100$, est
bien le produit des quatre racines, $5 \times 20$. La règle générale du traité
tient donc, exactement, sur l'exemple que le traité a construit pour montrer
qu'elle ne tient pas — à la seule condition de compter les côtés feints comme
des racines.\footnote{Calcul de cette lecture ; le feuillet écrit en propres
termes « $x$ vero fit summa duorum laterum de quibus est applicabilis », et cette
phrase est fausse sur ses propres nombres. C'est la démonstration la plus nette
que le volume fournisse de ce qui manque à sa doctrine, et elle ne demande
aucune connaissance que le feuillet n'ait pas : il suffit de comparer deux
nombres qu'il écrit tous les deux.}

Le chapitre se clôt sur trois remarques justes, qu'il faut rapporter parce
qu'elles montrent que la doctrine est cohérente dans ses limites. D'abord :
lorsque la question proposée ne peut se résoudre par les formules précédentes,
ou n'y tombe pas, « il deviendra manifeste que quelque côté feint a été assumé
dans sa composition, et ce côté sera enfin découvert ». C'est une méthode de
découverte par l'échec, et elle est correcte : c'est bien ainsi, faute de
critère direct, que la doctrine détecte les racines complexes. Ensuite : dans
les constitutions où $x$ est la différence des côtés, aucune équation anormale ne
peut naître, « car la somme et la différence étant données, le cas de la
question est toujours possible ». C'est l'observation, déjà faite au degré deux,
que le discriminant $(B+C)^2$ est toujours positif. Enfin : ce qui a été dit des
côtés feints pour le degré quatre vaut de même pour le degré trois.

La dernière phrase du traité renvoie les déterminations des quarto-quadratiques
à un traité compagnon qui n'est pas dans le volume, et s'arrête : « de iis plura
non dicemus ». Le reste de la page est blanc.

\section*{Le dossier de l'art de voler : ce qui est relié là}

\pagerange{26}{38}
Après une lacune de feuillets blancs, un tout autre dossier commence, sans
transition, sans titre général et sans rapport d'aucune sorte avec ce qui
précède. Il compte cinq pièces écrites et deux dessins.

Il ouvre sur un feuillet qui ne porte que deux mots au milieu de la page, en
capitales et minuscules droites : « Ars Volandi ». La page suivante est la page
de titre d'un livre imprimé, recopiée à la main, centrée, et dont voici le texte :

\begin{quote}
« De Arte volandi. Cuius ope, quivis homo, sine periculo, facilius, quam ullum
volucrem, quocumque lubet, brevissimo tempore ferri potest. Authore Friderico
Hermanno Flaydero, Poëtâ, professore et Bibliothecario Tubingæ. Cum eiusdem
operum affectorum et maximam partem perfectorum Indice. Typis Theodorici
Werlini, anno 1627. in 12\textsuperscript{o}. »\footnote{« De l'art de voler, par
le secours duquel tout homme peut, sans danger, plus aisément qu'aucun oiseau,
se porter où il lui plaît en très peu de temps. Par Friedrich Hermann Flayder,
poète, professeur et bibliothécaire à Tübingue. Avec l'index de ses ouvrages
entrepris et pour la plus grande part achevés. Des presses de Dietrich Werlin,
l'an 1627, in-12. » C'est la seule date de tout le volume, et elle est celle du
livre copié, non celle de la copie. Il s'ensuit seulement — et c'est l'unique
inférence chronologique que ces feuillets autorisent — que ce feuillet-là est
postérieur à 1627. L'inférence est de cette lecture. Aucune édition de Flayder
n'a été consultée, et cette lecture ne dit donc rien de la fidélité de la copie
au titre imprimé, ni du contenu du livre.}
\end{quote}

Suivent, au dos de ce feuillet et sur le suivant, des notes et extraits latins
sur le vol, sans titre, d'une cursive fine et très rapide, qui s'interrompent au
milieu d'une page. Puis, après deux vues blanches, la légende française des
lettres d'une figure, la figure elle-même — la machine dessinée en forme de
dragon —, le traité italien, et une seconde figure de la même machine vue de
face. Aucune des cinq pièces ne porte de date, aucune n'est signée.

\section*{Les notes latines : un dossier de témoignages}

\pagerange{28}{29}
Ces deux pages ne contiennent aucune mathématique et aucune expérience : ce sont
les notes de lecture d'un érudit qui rassemble, pour un livre à faire, tout ce
qu'il a trouvé sur des hommes qui se sont élevés en l'air. Elles valent d'être
inventoriées, parce que l'inventaire dit la nature de la pièce.

\begin{itemize}
  \item Les anciens ont appris à naviguer des oiseaux aquatiques et surtout du
    cygne, « dont le cou représente la proue, le ventre la carène, la queue la
    poupe, les ailes les voiles, les pattes les rames ». C'est un argument
    d'analogie et il ouvre le dossier : l'art imite l'animal.
  \item Le char à voiles du prince Maurice, attribué à « Stephinus », qui, sans
    le secours d'aucun cheval, « par la seule force des vents, voiles déployées,
    est emporté sur un rivage égal avec une telle rapidité qu'il laisse loin
    derrière lui les navires mêmes courant en pleine mer par le vent le plus
    favorable, et parcourt en l'espace de peu d'heures vingt ou trente milles
    d'Allemagne ». Et la remarque qui est le vrai argument de toute la page :
    « si ce char n'était pas encore inventé, qui croirait celui qui le
    raconterait ? »\footnote{« Stephinus » est Simon Stevin ; la transcription le
    note. L'argument — l'incrédulité n'est pas une preuve d'impossibilité — est
    exactement celui que le traité italien reprendra, et c'est le seul lien
    d'idée entre les deux pièces du dossier, sans qu'aucune ne renvoie à
    l'autre. Vingt à trente milles allemands font, selon la valeur retenue pour
    le mille, de cent cinquante à deux cent cinquante kilomètres ; l'ordre de
    grandeur est celui d'une exagération, et cette lecture ne l'évalue pas plus
    avant, le feuillet ne donnant ni le temps exact ni la distance.}
  \item Hadingus, roi des Danois, assiégeant Handuanus dans la ville de Duna,
    fit prendre des oiseaux du lieu et attacher à leurs plumes des amorces
    enflammées ; regagnant leurs nids, ils incendièrent la ville. « Saxo » et
    « Crantzius » sont donnés comme sources.
  \item Les pigeons du siège de Leyde, par lesquels la ville fut délivrée, et
    dont on rafraîchit encore la mémoire dans la ville même.
  \item Scaliger sur la colombe d'Archytas, avec le détail mécanique : la
    matière tirée de la moelle de jonc, revêtue des vessies ou pellicules dont
    se servent les batteurs d'or, enveloppée de nerfs, « où un demi-cercle, ayant
    poussé une roue, communiquera le mouvement aux autres, par lesquelles les
    ailes seront agitées ». C'est la seule description de mécanisme des deux
    pages, et elle est cohérente : un ressort de rappel entraînant un train de
    roues qui bat les ailes.\footnote{Le renvoi est donné « exercitationum
    Incardanum 326.\textsuperscript{a} », c'est-à-dire l'\emph{Exercitatio} 326
    des \emph{Exercitationes} de Scaliger contre Cardan. Aucune de ces deux
    œuvres n'a été consultée, et cette lecture ne vérifie pas la citation. Deux
    mots du passage, « facillimè » et « præstiterim », et un troisième,
    « neruulis » — lequel porte trois jambages initiaux et se lirait
    littéralement « meruulis » — sont donnés comme lectures incertaines par la
    transcription.}
  \item Le « Compendium miraculorum » d'Alvaro Gutiérrez de Torres, de Tolède,
    dédié à don Alonso de Fonseca, et le moine volant qu'il rapporte : « Elnero
    ex malmabiriâ », qui fut le plus docte de tous et prédit une comète bien
    avant son lever, et qui, dans sa jeunesse, s'attacha des plumes aux mains,
    croyant avoir trouvé l'art de voler « à l'instar de Dédale, dont il tenait
    la fable pour vraie » ; il s'élança d'une tour, parcourut plus d'un stade en
    volant, « c'est-à-dire plus de cent vingt-cinq pas », et fut jeté à terre,
    les jambes brisées, menant depuis une vie misérable ; il disait que l'unique
    cause de sa chute avait été d'avoir oublié une queue à la partie postérieure
    du corps.\footnote{Il s'agit du moine de Malmesbury, comme le note la
    transcription, laquelle donne « Elnero » comme lecture incertaine. Sur le
    fond, la cause que le récit lui fait alléguer est la bonne : sans surface
    stabilisatrice à l'arrière, un planeur n'a pas d'équilibre en tangage.
    C'est, par un détour de huit siècles, la fonction de l'empennage — et c'est
    la lettre D de la figure du dossier, « la queue laquelle fait tourner la
    machine du costé qu'on veut en haut et bas ». Le rapprochement est de cette
    lecture ; aucune des deux pièces ne renvoie à l'autre.}
  \item Ernst Burggrav, « médecin de ce siècle tout à fait admirable », et, dans
    son Planosphærium, un chantre de Nuremberg « qui, élevé dans l'air par la
    rame de deux ailes, s'envolait et redescendait à l'instar d'un oiseau », mais
    qui, par une imprudence — des roulettes fixées aux ailes s'embarrassaient ou
    ne s'appliquaient pas comme il fallait —, tomba et se brisa bras et jambes.
    Un cas semblable est rapporté de Paris, « de la mémoire de nos parents ».
  \item Johann Sturm, sur un homme qui se laissa descendre, ailé, de la tour
    Saint-Marc à Venise.
  \item Une liste des sens par lesquels les bêtes surpassent l'homme : « le
    vautour par l'odorat, le renard par l'ouïe, les chiens par l'un et l'autre,
    la poule par le goût, l'aigle par la vue, les limaces et les huîtres par le
    toucher, le lièvre à la course, le poisson à la nage, l'oiseau au vol ».
  \item Une note sur Flayder, et deux projets de livre de l'auteur des notes
    lui-même : un ouvrage sur les Bacchanales et leurs calamités, et un
    « Tractatus de villis Romanorum » où tout ce qui concerne les villas
    romaines serait recensé d'après les auteurs anciens, « mais aussi leurs
    ruines et leurs monuments subsistant encore aujourd'hui en Italie, décrits
    d'après de nouveaux itinéraires ».
\end{itemize}

Ces deux derniers points établissent la nature de la pièce mieux que tout le
reste : celui qui écrit ces notes tient un carnet de travail, il y mêle ses
lectures sur le vol et ses propres projets d'édition, et le dossier sur le vol
n'est qu'un de ses chantiers.\footnote{Une phrase de ces pages n'a pas été
rendue ici parce qu'elle ne se construit pas : « aquila plus homine pendet
uncijs quinque, quam in homine ponderis referri a philosophis memoratur » —
« l'aigle pèse cinq onces de plus que l'homme, qu'il n'est rapporté par les
philosophes de poids dans l'homme ». Le texte est altéré et cette lecture ne le
restitue pas ; il y est question d'un rapport de poids entre l'aigle et l'homme,
et rien de plus n'en est sûr. Trois autres endroits restent en suspens :
« præcurriōnem », que la transcription lit lettre à lettre
« p\textsuperscript{r}cuuriōnem » sans qu'aucun mot latin s'y ajuste ;
« philanthero », dans la note sur Flayder, également non résolu ; et la phrase
sur les Bacchanales, qui s'arrête sans verbe principal au bas d'une page.}

\section*{La machine, et la légende de ses lettres}

\pagerange{32}{37}
Deux dessins à la plume, au trait, sans lavis, montrent la même machine. Le
premier, de profil, est lettré ; le second, vu de face et de plus près, ne l'est
pas. Les figures ne sont pas reproduites ici, et les transcriptions les
décrivent. La machine a la forme d'un dragon : tête cornue, gueule ouverte,
corps écaillé en forme de tambour cerclé de bandes cousues, deux pattes
griffues, une queue s'achevant en large éventail palmé ; sur le dos, une calotte
bombée et quatre longues ailes étroites dressées deux par deux, entre lesquelles
deux cylindres annelés comme des vis ou des ressorts ; sous le corps, une grande
aile ; près de la tête, deux petites ailes palmées.

La page qui fait face au premier dessin en donne la légende, en français, d'une
autre main. Elle est brève et elle est la seule description du fonctionnement que
le volume contienne :

\begin{itemize}
  \item \textbf{A} — « aisle maistresse à laquelle de l'autre costé du dragon
    une autre est opposée, elles seruent toutes deux pour le pousser en auant et
    aussy pour le soustenir ». Deux grandes ailes, donc, chargées à la fois de
    la propulsion et de la portance.
  \item \textbf{B, B, B, B} — « quatre aisles qui se meuuent du haut en bas, et
    ne font autre chose que soustenir ». Quatre ailes portantes seulement.
  \item \textbf{C, C} — « deux petites aisles qui ne font que pousser en auant
    auec les deux maistresses ». Deux ailes propulsives seulement.
  \item \textbf{D} — « la queue laquelle fait tourner la machine du costé qu'on
    veut en haut et bas ». Une surface de gouverne, agissant en tangage et en
    lacet.
  \item \textbf{E} — « Couuerture mobile laquelle quand on veut s'eslargit
    circulairement sur la machine et la soustient assez pour l'empescher de
    precipiter ». Une voilure déployable au-dessus de l'appareil, destinée à
    freiner la chute.
\end{itemize}

Puis deux phrases sur la commande : « Un seul mouuement regle toutes les choses
encores qu'elles se meuuent en diuers temps », et « Toutes les aisles en se
leuant se ployent en elles mesmes et quand elles s'abbaissent elles
s'eslargissent, et le tout comme Je viens de dire est reglé par un seul
mouuement ».

Trois choses sont à dire sur cette légende, et la première est qu'elle contient
l'idée la plus juste de tout le dossier. Les ailes se replient à la montée et
s'élargissent à la descente : c'est exactement le principe du vol battu. Une
aile qui se referme en remontant présente moins de surface et engendre donc
moins de force vers le bas, tandis qu'ouverte en descendant elle engendre la
force maximale vers le haut ; le bilan sur un cycle est une force nette
ascendante, alors qu'une surface rigide battant symétriquement ne donnerait rien.
C'est ce que font les oiseaux, et la légende le décrit en une phrase.

La deuxième est que la pièce E est, en effet, un parachute, et que le principe en
est bon : une voilure déployée au-dessus d'une masse qui tombe augmente la
surface offerte au fluide et diminue d'autant la vitesse de chute, laquelle
varie comme l'inverse de la racine carrée de la surface. C'est la seule pièce de
la machine dont la fonction annoncée soit atteinte par le moyen indiqué, à ceci
près que la légende ne dit rien de la dimension nécessaire, qui est
grande.\footnote{Pour une masse de l'ordre de cent cinquante kilogrammes et une
vitesse de descente tolérable de cinq mètres par seconde, il faut une surface de
l'ordre de cent mètres carrés, soit un disque d'une dizaine de mètres de
diamètre : la « couuerture mobile » dessinée sur le dos du dragon n'en approche
pas. L'estimation est de cette lecture, et repose sur la traînée d'une voilure
en écoulement établi ; le feuillet ne donne aucune dimension.}

La troisième est que « un seul mouvement règle toutes les choses » décrit une
liaison à un seul degré de liberté, un train unique commandant les six ailes, la
queue et la couverture avec des phases différentes. C'est mécaniquement
concevable, et c'est peut-être la fonction des deux cylindres annelés que le
dessin place entre les ailes. Ni la légende ni le traité italien ne décrivent ce
mécanisme, et l'auteur du traité dit expressément qu'il ne le montrera
pas.\footnote{« Il dissegno della machina jnteriore, non mj par bene dj douerlo
mostrare » — « le dessin de la machine intérieure, il ne me paraît pas bon de
devoir le montrer » (vue 36 R). Le dessin intérieur n'est donc pas dans le
volume, et il n'y a rien à en dire.}

\section*{Le traité italien : le poids de l'air, et la flottaison}

\pagerange{33}{34}
Le traité italien commence par nommer son adversaire, qui est une opinion :
voler passe pour impossible au point que « le vulgaire, quand il veut affirmer
qu'une chose ne peut se faire, dit : il est aussi impossible de faire cela qu'il
est impossible de voler ». Puis l'auteur dit ce qu'il a fait, et c'est le seul
endroit du volume où quelqu'un parle de lui-même :

\begin{quote}
« Io però sino dallj miej primj annj applicaj l'animo a questa inuentione
parendomj strano che l'ingegno humano non fusse habile col suo sapere d'arriuare
a quelle cose che naturalj sono state concesse ad altrj animalj a noj jnferiorj
[…] sono gia decorsj x annj che da principij assaj communj mediante pero l'aiuto
diuino, ho ritrouato quello che qui sotto dimostrerò. »\footnote{« Cependant,
depuis mes premières années, j'ai appliqué mon esprit à cette invention,
trouvant étrange que l'esprit humain ne fût pas capable, par son savoir,
d'atteindre aux choses qui ont été naturellement accordées à d'autres animaux
inférieurs à nous […] voici déjà dix ans écoulés que, à partir de principes
assez communs, moyennant toutefois l'aide divine, j'ai trouvé ce que je
démontrerai ci-dessous. » Vue 33 R. L'auteur ajoute un argument
d'analogie qui est le même que celui du cygne des notes latines : la nature ne
nous a pas accordé d'habiter dans les eaux, et pourtant par l'étude nous
apprenons à nager. Il ne se nomme pas, et aucun feuillet ne le nomme ; cette
lecture ne le nomme donc pas.}
\end{quote}

Il raconte ensuite ce qu'il tient d'un tiers sur l'essai de Sagredo, et c'est un
récit précieux parce qu'il est quantitatif : Sagredo aurait pris un faucon,
l'aurait plumé, puis « observé diligemment la proportion que celles-ci — les
plumes — tenaient au corps, tant en poids qu'en mesure », et d'après cette
proportion se serait fait des ailes qu'il attachait à ses bras, « avec tant
d'artifice qu'en se jetant de quelque hauteur il arrivait à terre sans se faire
aucun mal, et s'écartait même de plusieurs brasses du pied de la hauteur ». Le
narrateur ajoute qu'il a lui-même fait beaucoup de tentatives et que « toute la
peine fut prise en vain ».\footnote{L'auteur nomme « Gio. Fran. Sagredo, noble
vénitien, très noble mathématicien », et « il Sig.\textsuperscript{r} Galileo
Galilej, noble florentin et autre Archimède de nos temps ». La transcription donne
l'initiale du second prénom comme incertaine, le nom étant écrit « Gio. Ra~n »
avec un trait de suspension. La dernière phrase du récit, « ma giu non poteua
tramontare per niuna via », n'est pas rendue ici parce que son sens n'est pas
sûr : elle dit que Sagredo « ne pouvait tramontare vers le bas par aucune
voie », ce qui s'entendrait aussi bien d'une impossibilité de descendre que
d'une impossibilité de se poser autrement ; cette lecture ne tranche pas.}

Le raisonnement quantitatif de mise en train de la flottaison occupe ensuite
trois pages, et il se suit pas à pas.

\subsection*{Le poids de l'air}

L'auteur pose qu'Aristote a écrit que tous les éléments ont de la pesanteur, et
demande dans quelle proportion l'air pèse à l'eau. « Selon Aristote ce serait
comme un à dix, mais selon le calcul de M. Galilée ce serait, et il est plus
vrai, comme un à quatre cents. »

Il faut dire ce qui est vrai. À la température et à la pression ordinaires,
l'eau pèse environ mille kilogrammes par mètre cube et l'air environ un et deux
dixièmes : le rapport est donc d'environ un à huit cents, et il varie de un à
sept cent soixante-dix à zéro degré à un à huit cent trente à vingt degrés. Le
chiffre d'Aristote est faux d'un facteur quatre-vingts. Le chiffre que le
feuillet attribue à Galilée est faux d'un facteur deux, dans le même sens :
l'air y est deux fois plus lourd qu'il n'est. Le jugement du feuillet — le second
est « plus vrai » — est donc entièrement fondé, et l'écart qui reste est de ceux
qu'une mesure ancienne du poids de l'air pouvait laisser.\footnote{Aucune
édition de Galilée n'a été consultée, et cette lecture ne dit donc rien de ce
que Galilée a écrit : elle rapporte ce que le feuillet lui attribue, et elle
compare ce chiffre à la valeur admise aujourd'hui. Le chiffre importe pour la
suite, parce que toute la conclusion sur la flottaison en dépend, et qu'elle
tombe avec lui.}

\subsection*{Archimède, et les pesanteurs spécifiques}

Le traité cite ensuite, comme « fondement de tout mon discours », une proposition
d'Archimède sur les corps qui flottent :

\begin{quote}
« Le cose piu grauj poste nell'humido discendono tanto che vanno al fundo e sj
fanno tanto piu leggieri nell'humido, quanto é il peso dell'humido che habbia
tanta grandizza quanta é la grandezza dj dette cose. »\footnote{« Les choses plus
pesantes placées dans l'humide descendent jusqu'à ce qu'elles aillent au fond,
et se font dans l'humide d'autant plus légères qu'est le poids de l'humide qui
aurait autant de grandeur qu'est la grandeur desdites choses. » Vue 34 L. Les
quatre lignes de la citation sont précédées, ligne à ligne, d'un même signe
bouclé dans la marge, qui est la marque de citation de cette main.}
\end{quote}

L'énoncé est exact : un corps plus dense que le fluide coule, et son poids y est
diminué du poids du fluide qu'il déplace. C'est ce que nous appelons la poussée
d'Archimède, et le traité la tient dans la forme juste, celle du poids du fluide
de volume égal.\footnote{Le nom de poussée d'Archimède est postérieur au
feuillet, l'énoncé lui est attribué par le feuillet même, et la coïncidence est
complète. C'est le seul énoncé de tout le dossier du vol qui soit à la fois cité
d'un auteur et parfaitement correct.}

Viennent alors cinq nombres, les pesanteurs spécifiques rapportées à l'eau :
« l'or tient proportion à l'eau comme 20 à l'unité, le mercure comme
$13\tfrac12$, le plomb comme 11, l'argent comme 10 et le cuivre comme 9 ». Voici
ce que valent ces cinq nombres aujourd'hui, à la température ordinaire : or
$19{,}3$ ; mercure $13{,}5$ ; plomb $11{,}3$ ; argent $10{,}5$ ; cuivre $9{,}0$.
L'ordre est juste, et les cinq valeurs sont exactes à cinq pour cent près, le
mercure et le cuivre à moins d'un pour cent. C'est un bon tableau, et il faut le
dire.\footnote{Comparaison de cette lecture. Écarts : or $+3{,}6\%$, mercure
$-0{,}2\%$, plomb $-3\%$, argent $-4{,}7\%$, cuivre $+0{,}4\%$. Le feuillet ne
dit pas d'où il les tient.}

Sur quoi le traité tire une conséquence qui demande une correction. Il écrit que
ces métaux arriveront au fond dans cet ordre, l'or le premier et le cuivre le
dernier, « et cela selon leur pesanteur », et qu'ils feront de même dans l'air,
mais plus vite, l'air ayant moins de pesanteur que l'eau. Ce qui est vrai est
ceci : dans un fluide résistant, la vitesse limite d'un corps croît avec l'écart
entre sa densité et celle du fluide, mais elle dépend aussi de sa taille et de
sa forme ; l'ordre annoncé ne vaut donc que pour des corps de même grandeur et
de même figure, condition que le feuillet ne pose pas. Et dans le vide, tous les
corps tombent de la même manière, ce que le cadre du traité ne peut pas
atteindre, puisqu'il rattache la vitesse de chute à la seule pesanteur
spécifique.\footnote{Départ de cette lecture par rapport au feuillet. Le feuillet
a raison sur le second point qu'il avance — les corps descendent plus aisément
dans l'air que dans l'eau parce que « le milieu est plus léger » — et cette
raison est la bonne : c'est la résistance du milieu, non la poussée seule, qui
gouverne. Sa doctrine générale est celle d'un mouvement dont la vitesse est
réglée par la pesanteur spécifique, laquelle n'est vraie que là où la résistance
domine.}

\subsection*{La flottaison dans l'air, et le corps plus léger que l'air}

Le traité arrive alors à son argument central, et il est bien conduit. Si l'on
trouvait une matière qui fût à l'eau comme 1 à 400, elle serait « en pesanteur
égale à l'air, c'est-à-dire qu'elle ne monterait pas dans l'air et n'y
descendrait pas non plus, à moins d'y être poussée violemment en haut ou en
bas ». Et si l'on en trouvait une qui fût comme 1 à 500, « celle-là, d'elle-même
et sans aucune violence, monterait dans l'air jusqu'à la sphère du feu, où elle
s'arrêterait ». Alors, conclut-il, « nous pourrions très facilement former un
instrument et voler avec lui comme font les oiseaux ».

Le raisonnement est juste et la conclusion est fausse, pour deux raisons
distinctes qu'il faut séparer.

La première est le nombre. La flottaison neutre a lieu quand la densité moyenne
du corps égale celle du fluide : cela, le traité l'a exactement. Mais la densité
de l'air ne vaut pas le quatre-centième de celle de l'eau, elle vaut environ le
huit-centième. Une matière au quatre-centième est donc deux fois plus lourde que
l'air et tombe. C'est le chiffre emprunté qui fait tomber la conclusion, non le
raisonnement : avec la bonne valeur, il faudrait une matière deux fois plus
légère encore que celle que le traité réclame.\footnote{Conséquence tirée par
cette lecture des nombres du feuillet. Elle mérite d'être notée parce qu'elle
montre que l'auteur est plus près de la vérité qu'il ne le croit sur le principe
et plus loin qu'il ne le croit sur la faisabilité.}

La seconde est la « sphère du feu ». Il n'y en a pas. Un corps moins dense que
l'air monte, mais il ne monte pas indéfiniment ni jusqu'à une région assignée du
monde : la densité de l'air décroît avec l'altitude, et le corps s'arrête là où
la densité de l'air environnant est descendue jusqu'à la sienne. Il y a bien une
altitude d'arrêt, ce que le traité pressent en écrivant « où elle s'arrêterait »,
et elle est déterminée par la décroissance de la densité, non par la limite d'un
élément.\footnote{Départ de cette lecture. La cosmologie des quatre éléments et
de leurs sphères est celle du feuillet, qui du reste en doute pour le feu : « et
du quatrième, s'il y en a un, c'est-à-dire du feu, il ne nous importe pas d'en
dire davantage ». Le nom moderne de ce qui remplace la sphère du feu est
l'équilibre hydrostatique de l'atmosphère, et le fait qu'un ballon possède une
altitude d'équilibre en est la conséquence ; les deux sont postérieurs au
feuillet.}

L'auteur ne s'arrête d'ailleurs pas à cette voie : il constate aussitôt qu'aucune
telle matière n'est à notre portée, et que les oiseaux n'en ont pas non plus. Il
observe que les plumes, « très légères » soient-elles, tombent quand on les
laisse aller d'en haut, même les plus menues ; et que si l'on jetait un oiseau
d'une hauteur les ailes ouvertes mais liées, il ne se maintiendrait pas en l'air
et tomberait avec une percussion notable, « parce que son soutènement ne provient
d'autre chose que de la percussion des ailes dans l'air ». C'est là le pivot du
traité, et l'observation est juste : la portance d'une aile battante n'est pas
une flottaison, c'est une force de réaction exercée par l'air.

\section*{La percussion}

\pagerange{35}{36}
L'auteur ouvre une digression sur la percussion, « qui est à mon avis l'une des
plus grandes choses qui soient dans toute la mécanique », et il donne la raison
de l'admirer :

\begin{quote}
« vediamo superare la forza dj 1000 e 2000 libre dj resistenza con vn martello
che pesj solo x o xii libre. »\footnote{« Nous voyons vaincre la force de 1000 et
2000 livres de résistance avec un marteau qui ne pèse que 10 ou 12 livres. »
Vue 35 L. Il ajoute que la chose paraît d'autant plus merveilleuse que
« personne jusqu'ici, que je sache, n'a écrit là-dessus ». Cette remarque est un
énoncé sur l'état de son information, et cette lecture ne la vérifie pas : elle
demanderait de consulter ce qui avait été imprimé, ce que la règle de cette
édition interdit.}
\end{quote}

L'observation est juste et même modeste. Ce qu'un marteau transmet n'est pas son
poids mais sa quantité de mouvement, et la force moyenne du choc est cette
quantité divisée par la durée du contact, laquelle est très brève : c'est ce que
nous appelons l'impulsion,
\[ F\,\Delta t = m\,\Delta v. \]
Un marteau de cinq kilogrammes arrivant à dix mètres par seconde et arrêté en un
millième de seconde exerce une force moyenne de cinquante mille newtons, soit
plus de cinq tonnes-force, c'est-à-dire beaucoup plus que les deux mille livres
du feuillet.\footnote{Calcul de cette lecture, avec ses hypothèses : dix ou douze
livres font environ cinq kilogrammes, la durée du contact d'un choc de métal sur
métal est de l'ordre du millième de seconde. Le nom d'impulsion et la relation
écrite sont postérieurs au feuillet, qui n'a ni l'un ni l'autre ; ce qu'il a,
c'est le fait et l'étonnement juste devant le fait.}

Il divise ensuite la percussion en trois espèces, et la division est bonne :

\begin{itemize}
  \item \textbf{ferma}, arrêtée : quand on frappe du marteau sur un solide
    immobile posé sur un plan, comme sur une enclume ou sur une pierre ;
  \item \textbf{incontrata}, rencontrée : quand on frappe sur un solide mobile
    qui vient à la rencontre du coup, « comme serait un marteau contre un autre,
    ou une épée contre une autre ; et celle-ci est la plus violente de toutes » ;
  \item \textbf{sfuggita}, fuie : quand on frappe sur un solide mobile qui fuit
    le coup, « et celle-ci est la plus faible de toutes ».
\end{itemize}

Et il en donne la loi, sur un exemple qui est exactement le bon :

\begin{quote}
« se il moto del vento sara eguale al moto della percossa, questa sera di niun
valore e sara à guisa dj doj caualljerj che vno carichj l'altro con la lancia,
che se la velocita dell'vno sara eguale a quella dell'altro, il percuttore non
potra maj ferire quello il fuggitiuo benche hauesse la punta della lancia adosso
il corpo dell'jnimico. »\footnote{« Si le mouvement du vent est égal au mouvement
de la percussion, celle-ci sera de nulle valeur, et ce sera à la manière de deux
cavaliers dont l'un charge l'autre avec la lance : si la vitesse de l'un est
égale à celle de l'autre, celui qui frappe ne pourra jamais blesser le fuyard,
quand même il aurait la pointe de la lance contre le corps de l'ennemi. »
Vue 36 L.}
\end{quote}

C'est juste, et c'est plus juste que l'auteur ne peut le dire. L'impulsion
échangée dans un choc ne dépend que de la vitesse \emph{relative} des deux
corps ; si cette vitesse relative est nulle, il n'y a pas de choc, quelle que
soit la vitesse commune. Et la vitesse relative est la même dans tous les
repères en translation uniforme, ce qui est la raison de fond pour laquelle la
division en trois espèces est une division de la vitesse relative et non de trois
phénomènes différents : une même percussion est ferme, rencontrée ou fuie selon
ce que fait le corps frappé.\footnote{L'invariance de la vitesse relative par
changement de repère galiléen est postérieure au feuillet, qui n'a pas le
concept de repère. Ce qu'il a, et qui est remarquable, c'est l'exemple des deux
cavaliers, qui isole exactement la bonne quantité.}

Suivent les percussions sur les corps rares, l'eau et l'air, avec quatre
observations qu'il vaut de reprendre une à une.

\begin{itemize}
  \item Une épée très forte frappant l'eau à plat se brise « le plus souvent en
    morceaux », alors qu'elle avait tenu à d'autres épreuves ; et du bois tombant
    de quelque hauteur se brise dans l'eau, qui aurait suffi, là où il fut rompu,
    à soulever un grand poids. C'est juste : l'eau, à l'échelle de temps d'un
    choc, ne s'écarte pas assez vite pour se dérober, et la lame doit mettre en
    mouvement une masse d'eau considérable en très peu de temps ; la réaction qui
    en résulte excède la résistance de la lame.
  \item Un boulet d'arquebuse tiré obliquement sur l'eau d'un large fleuve ne
    s'y enfonce pas mais, la frappant par incidence, « se porte à l'autre rive en
    sautelant, et ira parfois deux ou trois fois plus loin qu'il n'est naturel
    dans l'air ». Le ricochet est réel et l'explication est bonne. La comparaison
    est défendable si on l'entend, comme le texte y invite, d'un tir à faible
    élévation : sur terre le boulet s'enfonce au premier contact, sur l'eau il
    est renvoyé et continue. Il faut seulement ajouter que chaque rebond lui
    enlève de l'énergie, en sorte que la distance totale ne peut pas dépasser la
    portée qu'aurait, sans obstacle, la même charge.\footnote{Précision de cette
    lecture ; le feuillet ne dit rien de la perte à chaque rebond, et sa phrase,
    prise absolument, suggérerait un gain de portée, qui est impossible.}
  \item L'air. « La percussion ferme dans l'air est celle qui se fait par temps
    tranquille sur quelque solide violemment mû », et l'auteur rapporte avoir vu
    souvent une planchette ou autre chose mince, tournée avec furie dans l'air,
    se briser au milieu de sa course, elle qui sans cela aurait suffi à soulever
    un poids sept ou huit fois plus grand que le sien.
  \item Le roseau. Placé horizontalement dans un mur, sa cime s'incline vers le
    centre par sa propre pesanteur ; mais tenu à la main verticalement, pointe en
    haut ou en bas, et mû violemment de la droite ou de la gauche, sa cime
    « s'inclinera vers le commencement du mouvement, c'est-à-dire vers le point
    de son zénith ou de son nadir », et se brisera si le mouvement est assez
    violent. L'auteur conclut que « cela ne provient d'autre chose que du
    mouvement auquel l'air résiste », et que la percussion surpasse de beaucoup
    la pesanteur du roseau.
\end{itemize}

Sur ce dernier point il faut une correction, et elle porte sur la cause, non sur
le fait. Le fait est exact : le roseau fléchit en arrière et peut rompre, et la
charge transversale du mouvement excède largement celle de son propre poids,
comme le montre la comparaison des deux situations — cette conclusion du feuillet
est bonne. Mais la cause principale n'est pas la résistance de l'air : c'est
l'inertie du roseau lui-même. Sa masse est répartie sur toute sa longueur, et
chaque élément doit être accéléré ; la pointe, la plus éloignée de la main,
demande la plus grande accélération et reste en arrière parce qu'il faut une
force pour l'y porter. La résistance de l'air s'y ajoute et n'y suffit pas ; pour
un mouvement rapide, l'inertie domine.\footnote{Départ de cette lecture. Le
feuillet nomme tout de même le mouvement — « dal moto alquale l'aria resiste »,
du mouvement auquel l'air résiste — et non l'air seul, en sorte que la faute est
de partage plutôt que d'aveuglement. La notion d'inertie d'un corps étendu et le
calcul du moment fléchissant d'une poutre en rotation sont postérieurs au
feuillet.}

\section*{Le levier : on ne gagne pas ensemble la force et le temps}

\pagerange{36}{36}
Le traité passe alors à ce qu'il appelle « l'antipathie » du mouvement et de la
force, et il la démontre sur un levier. C'est la page la plus soignée du traité,
et son argument est juste.

\begin{quote}
« È impossibile come mecanicamente sj proua in vn jstesso tempo auanzar tempo e
forza, perche quanta maggior forza sj acquista all'incontra altretanto tempo sj
perde. »\footnote{« Il est impossible, comme on le prouve mécaniquement, de
gagner en un même temps le temps et la force, parce que d'autant plus de force
on acquiert, d'autant de temps en revanche on perd. » Vue 36 L.}
\end{quote}

La figure est un levier horizontal $AB$ appuyé en $C$, le point d'appui étant au
quart de la longueur du côté de $A$ ; sous $A$ pend un petit cercle contenant le
chiffre 4, sous $B$ un cercle contenant le chiffre 1 ; une droite oblique
descend de $D$, au-dessus de $A$, traverse le levier en $C$ et va jusqu'à $E$,
au-dessous de $B$ ; deux arcs pointillés joignent $A$ à $D$ et $B$ à $E$. Le
raisonnement du texte est : une livre de poids ou de force placée en $B$
soulèvera ou équilibrera quatre livres placées en $A$, parce que la distance
$CB$ est quadruple de $CA$ ; mais « à la force placée en $B$ il faudra quatre
fois plus de temps pour soulever le poids placé en $A$ », puisque, pour pousser
$A$ jusqu'en $D$, il faut que $B$ s'abaisse jusqu'en $E$, et que l'arc $BE$ est
quadruple de l'arc $AD$.

Tout cela est exact. Écrivons-le. L'équilibre du levier donne
\[ F_A \cdot CA = F_B \cdot CB, \qquad \text{d'où} \quad F_B = \tfrac14 F_A \ \text{si} \ CB = 4\,CA, \]
et la rigidité du levier donne, pour les déplacements des deux extrémités,
\[ \frac{\delta_A}{CA} = \frac{\delta_B}{CB}, \qquad \text{d'où} \quad \delta_B = 4\,\delta_A. \]
En multipliant, $F_B\,\delta_B = F_A\,\delta_A$ : la quantité conservée est le
produit de la force par le déplacement, c'est-à-dire le \emph{travail}. Ce qu'on
ne peut pas gagner ensemble, c'est la force et le déplacement ; et c'est le
principe que nous appelons aujourd'hui principe des travaux
virtuels.\footnote{Le nom et la notion de travail sont très postérieurs au
feuillet. L'énoncé du feuillet et celui-ci coïncident, mais non littéralement, et
c'est le point suivant.}

Car il y a, dans la lettre du feuillet, une confusion qu'il faut lever : il dit
« temps » là où sa figure montre des \emph{chemins}. Sur un levier rigide, les
deux extrémités se meuvent en même temps : ce qui est quadruple n'est pas la
durée, c'est l'arc parcouru, et donc la vitesse. L'énoncé du feuillet redevient
vrai dès qu'on lui fournit l'hypothèse qu'il n'écrit pas — que la main qui agit
en $B$ se meut à une vitesse donnée : alors, pour élever $A$ d'une hauteur
donnée, il faut bien quatre fois plus de temps, puisqu'il faut parcourir quatre
fois plus de chemin à vitesse égale. Mais l'hypothèse n'est pas posée, et la
figure ne montre qu'un rapport d'arcs.\footnote{Départ de cette lecture par
rapport au feuillet. C'est une correction de langage et non de fond : la
conclusion que le traité tire — on ne gagne pas ensemble la force et le temps —
est vraie sous l'hypothèse tacite, et la conclusion qu'il en fait servir pour sa
machine reste valable. Le traité écrit du reste aussi bien le cas réciproque, « si
le moment est en $A$, où il faudra quatre fois plus de force pour mouvoir
l'extrémité $B$, mais nous gagnerons quadruplement en temps », et là encore le
rapport est celui des chemins.}

De là le traité tire une conséquence sur Archimède, et le raisonnement est bon.
Il lit dans la vie de Marcellus par Plutarque qu'Archimède faisait des machines
qui, descendues des murs, saisissaient les galères romaines dans le port de
Syracuse et « avec grande promptitude les élevaient en haut », puis les
laissaient retomber dans l'eau pour les submerger, en sorte que les Romains
l'appelaient « Briareo geometra ». « Ce qui, si c'est vrai, veut dire qu'Archimède
avait trouvé le moyen de multiplier la force et le temps ensemble, chose qui non
seulement ne se sait pas aujourd'hui, mais est même tenue pour impossible. » Il
ajoute que l'autre épreuve rapportée du même Archimède, avoir tiré seul à terre
une très grosse caraque que beaucoup d'hommes n'avaient pu remuer, ne lui paraît
nullement impossible, « sachant très bien que la force peut se multiplier à
l'infini ».

Les deux jugements sont justes. Aucune machine ne multiplie à la fois la force
et la vitesse : si les grues d'Archimède multipliaient énormément la force, elles
ne pouvaient pas soulever vite, et l'un des deux termes du récit est exagéré. Et
la force, en revanche, se multiplie sans limite de principe par un rapport de
bras suffisant — la limite étant pratique, la vitesse décroissant dans la même
proportion, et les frottements et la résistance des pièces intervenant. Le
traité a exactement ce qu'il faut pour trancher, et il
tranche.\footnote{L'auteur ne consulte pas Plutarque autrement qu'en le citant,
et cette lecture ne l'a pas consulté du tout : ce qui est rapporté ici est ce
que le feuillet dit de Plutarque, et rien de plus. Le raisonnement, lui, ne
dépend pas de la source.}

\section*{Ce que la machine pouvait faire}

\pagerange{36}{37}
Le traité tire de son propre principe la limite de sa machine, et il le fait
honnêtement :

\begin{quote}
« l'istromento da me jnuentato (a mio parere) non sj potra far maggiore che
possj portare o tenir dentro piu dj doj huominj per che non solo per mouere le
alj vuole gran forza ma ancora vi vuole gran prestezza, le qualj doj cose come
ho mostrato fra dj loro hanno grandissima antipatia. »\footnote{« L'instrument
inventé par moi, à mon avis, ne pourra se faire plus grand qu'il ne puisse
porter ou tenir dedans plus de deux hommes, parce que non seulement pour mouvoir
les ailes il faut une grande force, mais il y faut encore une grande promptitude,
lesquelles deux choses, comme je l'ai montré, ont entre elles la plus grande
antipathie. » Vue 36 R.}
\end{quote}

Il ajoute avoir « pour ainsi dire trompé la nature » en se servant d'une
multiplicité de mouvements contraires l'un à l'autre, « lesquels ne sont pas de
peu de secours à soulever, ce qu'ils ne feraient pas s'ils étaient simples ». Et
il décrit sa machine : faite en forme de dragon, elle peut tenir jusqu'à deux
hommes, dont il suffit qu'un seul travaille tandis que l'autre se repose ; elle
peut cheminer de nuit avec l'aide de la boussole, portant de quoi manger et
boire pour quelques jours ; et — c'est sur cette phrase que le traité s'achève —
elle est faite de telle manière que si l'une des ailes venait à se rompre, elle
ne pourrait pour autant précipiter, mais tomberait doucement, en sorte que ceux
qui sont dedans « souffriraient très peu de lésion ».

Il faut dire ce qu'il en est, et le dire sans mépris, parce que le raisonnement
qui mène là est de bonne qualité.

\subsection*{L'échelle : pourquoi les ailes de Sagredo ne pouvaient pas marcher}

Le procédé prêté à Sagredo — mesurer sur un faucon la proportion des plumes au
corps, en poids et en mesure, et se faire des ailes dans la même proportion — ne
peut pas réussir, et la raison en est géométrique. Quand on agrandit un animal en
gardant ses proportions, sa surface croît comme le carré de sa dimension
linéaire et sa masse comme le cube : la charge par unité de surface d'aile croît
donc comme la dimension. Un faucon pèse un peu plus d'un demi-kilogramme, un
homme soixante-dix : le rapport des masses est de l'ordre de cent vingt, celui
des dimensions de l'ordre de cinq, et la charge alaire est multipliée par cinq.
Comme la vitesse nécessaire pour soutenir cette charge croît comme sa racine
carrée et la puissance comme le produit du poids par cette vitesse, la puissance
demandée est multipliée par plusieurs centaines, alors que la force des muscles
ne croît, au mieux, que comme leur masse, c'est-à-dire par cent vingt. La
proportion, qui était le principe même de la méthode, est précisément ce qui ne
se conserve pas.\footnote{Calcul d'ordres de grandeur de cette lecture, avec ses
hypothèses écrites. Les noms modernes en sont la charge alaire et la loi
carré-cube, tous deux postérieurs au feuillet ; l'argument, lui, ne demande que
de comparer un carré à un cube, et rien dans les connaissances que le feuillet
met en œuvre n'empêchait de le faire. Le feuillet ne le fait pas et se borne à
constater l'échec des essais.}

\subsection*{La puissance : ce que les nombres du feuillet permettent d'établir}

La machine ne pouvait pas voler, et on peut le montrer avec les nombres que le
feuillet lui-même fournit. Soutenir un poids en battant l'air, c'est communiquer
continuellement à l'air une quantité de mouvement vers le bas ; si $W$ est le
poids, $\rho$ la masse par unité de volume de l'air et $S$ la surface balayée par
les ailes, la vitesse qu'il faut donner à l'air est
\[ v = \sqrt{\frac{W}{2\rho S}}, \]
et la puissance à fournir est $P = W v$. Pour deux hommes et leur machine, soit
un poids de l'ordre de mille cinq cents newtons, et une surface balayée
généreusement estimée à vingt mètres carrés, on trouve $v \approx 5{,}6$ mètres
par seconde et $P \approx 8\,400$ watts. Un homme entraîné fournit de façon
soutenue deux à trois cents watts. Le déficit est d'un facteur de l'ordre de
trente.

Et le feuillet ne peut pas s'en tirer par son propre chiffre sur l'air. Avec la
densité qu'il lui donne — le quatre-centième de l'eau, soit deux et demi
kilogrammes par mètre cube au lieu d'un et deux dixièmes —, l'air est plus dense,
donc plus facile à frapper, et la puissance requise tombe à environ six mille
watts : encore vingt fois trop. Le calcul est le même quel que soit celui des
deux chiffres qu'on retienne.\footnote{Calcul et estimations de cette lecture,
avec toutes ses hypothèses : la relation employée est celle du bilan de quantité
de mouvement pour un disque sustentateur en vol stationnaire, la puissance
musculaire soutenue d'un homme est celle qu'on mesure aujourd'hui, et la surface
balayée est une majoration généreuse de ce que le dessin montre. Rien de tout
cela n'est sur le feuillet, et la conclusion est donc celle de cette lecture et
non une réfutation que l'auteur aurait pu écrire. Il faut ajouter, pour être
juste, que ce déficit de puissance est la raison pour laquelle le vol humain par
la seule force des muscles n'a été obtenu qu'au XX\textsuperscript{e} siècle, et
qu'il l'a été avec une aile fixe et une hélice, non avec des ailes battantes, au
moyen de structures d'une légèreté hors de portée de tout atelier ancien.}

Le fait remarquable est que l'auteur avait entre les mains le principe qui
condamne sa machine, et qu'il s'en est servi pour en limiter la taille sans aller
jusqu'à la conclusion. Son propre théorème — on ne gagne pas ensemble la force et
la vitesse — dit que le mouvement des ailes demande à la fois de la force et de
la promptitude, et que ces deux exigences se combattent. Il en tire : donc pas
plus de deux hommes. Ce qu'il fallait en tirer : donc pas même un. Ce qui lui
manquait pour faire le pas n'était pas la mécanique, c'était la mesure du poids
de l'air et l'idée de compter une puissance.

\subsection*{Les deux dernières promesses}

Deux affirmations de la fin ne sont pas soutenues par le feuillet et il faut le
dire. La première est que si une aile se rompt, la machine ne précipite pas mais
descend doucement : rien dans le texte ni dans les figures ne montre comment, et
une machine dissymétrique par la perte d'une aile n'a aucune raison de descendre
droit. La seule pièce qui pourrait le faire est la « couuerture mobile » E de la
légende française, dont la légende dit précisément qu'elle « soustient assez pour
l'empescher de precipiter » ; mais la légende et le traité sont de deux mains, et
ni l'un ni l'autre ne les rapproche. La seconde est que ceux qui sont dedans
souffriraient très peu de lésion : c'est une promesse, non une conséquence.

Il faut être juste avec ce texte. Son auteur ne promet pas un prodige : il
raisonne, il chiffre ce qu'il peut chiffrer, il emprunte le meilleur chiffre
qu'il connaisse sur le poids de l'air, il cite Archimède dans la forme juste,
il isole la bonne quantité dans un choc, il démontre correctement la loi du
levier et il s'en sert contre lui-même pour brider son propre projet. C'est le
travail d'un mécanicien, et il est faux par les endroits où, sur ce sujet, on
pouvait alors se tromper : le poids de l'air, et l'absence de toute notion de
puissance.

\section*{Ce que ces feuillets ne disent pas}

\pagerange{2}{38}
Il faut finir par les négations, parce que le titre du catalogue mène droit à
quatre erreurs et qu'un lecteur qui arrive par lui doit être averti.

\textbf{Il n'y a aucune lettre dans ce volume.} Le titre annonce « des opuscules
et des lettres de P. de Roberval, Huggens, Torricelli, Firmat et Fr. Herman
Flayder » : la correspondance de Roberval, de Huygens, de Torricelli et de Fermat
n'est dans aucun des deux lots, ni en original, ni en copie, ni par extrait. De
Flayder, il n'y a pas d'opuscule non plus : il y a une page de titre recopiée à
la main.

\textbf{Le traité d'algèbre n'est attribué à personne ici.} Il est anonyme sur le
feuillet : ni nom d'auteur, ni souscription, ni dédicace. Roberval y paraît
deux fois sous la forme « D. Rob. », aux vues 9 R et 10 R, comme l'auteur de
trois lemmes que le traité déclare emprunter, avec une note marginale en
français qui renvoie à « son traité de mechanique des plans inclinez ». C'est
une citation. Attribuer à un auteur un texte parce qu'il y est cité serait
retourner la preuve, et cette lecture ne le fait pas — pas plus qu'elle
n'attribue le traité à Viète, dont il écrit la notation, ou à quiconque d'autre.
Le catalogue range le volume sous le nom de Roberval ; les feuillets ne le
disent pas.

\textbf{Le traité italien n'est attribué à personne non plus.} Son auteur nomme
Sagredo et Galilée, il parle abondamment de lui-même — ses premières années, ses
essais vains, les dix ans écoulés depuis sa trouvaille, sa machine — et il ne se
nomme jamais. Aucun feuillet du dossier ne le nomme. Cette lecture ne le nomme
pas.

\textbf{Aucun feuillet ne porte de date.} La seule date du volume, 1627, est à
l'intérieur de la page de titre recopiée, et c'est celle du livre imprimé, non
celle de la copie. Le champ de datation du catalogue est vide et cette lecture
le laisse vide. Deux inférences sont possibles à partir du contenu, et elles sont
signalées comme inférences là où elles sont faites : le feuillet qui copie le
titre de Flayder est postérieur à 1627, et le traité italien est postérieur à
l'estimation du poids de l'air qu'il cite. Rien d'autre. Le traité d'algèbre ne
donne, pour tout repère, que son emprunt à Roberval et la notation de l'école de
Viète, et cette lecture n'en tire aucune année.

\textbf{Enfin, il n'y a pas un texte, mais deux dossiers.} La tentation, devant
un volume relié, est de chercher ce qui tient ses pièces ensemble : une même
main, un même projet, un fil. Il n'y en a pas ici. Un traité d'algèbre sur la
composition des équations et un dossier sur l'art de voler ont été cousus dans la
même reliure, avec une lacune de feuillets blancs entre eux, et rien dans l'un ne
fait signe vers l'autre. On peut relever, et cette lecture l'a fait, que les deux
dossiers partagent une façon de raisonner — l'énumération exhaustive des cas dans
le premier, la division de la percussion en trois espèces dans le second — et
qu'ils citent l'un et l'autre un mathématicien contemporain pour un lemme de
mécanique. C'est une ressemblance d'époque, non un lien. Prétendre autre chose
serait ajouter au volume ce que le relieur seul y a mis.

\end{document}
